% 用法: xelatex SL_gap_n1_O3a_phase_rigidity_proof.tex (两遍, -output-directory=build)
% 主题: O3a 完整证明 - 一维加权弦势垒族中 sign-consistent good root 的唯一性
% 来源: 用户提供 O3a_complete_proof_zh.pdf (2026-08-09), 本文为忠实转录;
%       证明包 SHA-256 7e49f975a0dcd3cdc639515d5f2d72a004a4c36322991c067ce0f492c0a9f366,
%       合并后 canonical Blueprint SHA-256 95cc5b2f7f8c23526642344a8fe1b6139611b73f679c3b90f2f37d9f654b303b.
% 审计: 2026-08-09 符号/数值审计 scripts/audit_o3a_pdf_part1..4.py 等, 全部通过 (详见第 8 节前的说明).
% 审计补充 (2026-08-09 续): 附录 lem:envseries 与表题注的表述修正 (F-208 sin 夹逼方向, F-209 arctan 项数与包络宽度), 证书数据不变.
\documentclass[UTF8]{ctexart}
\usepackage{amsmath, amssymb, amsthm}
\usepackage[margin=2.5cm]{geometry}
\usepackage[hyphens]{url}
\usepackage[hidelinks]{hyperref}
\usepackage{booktabs}
\usepackage{longtable}
\usepackage{enumitem}
\emergencystretch 3em

\newtheorem{theorem}{定理}[section]
\newtheorem{proposition}[theorem]{命题}
\newtheorem{lemma}[theorem]{引理}
\newtheorem{corollary}[theorem]{推论}
\newtheorem{remark}[theorem]{注}
\newtheorem{definition}[theorem]{定义}

\title{O3a 完整证明: 一维加权弦势垒族中 sign-consistent good root 的唯一性}
\author{BVE research 项目 (基于 Blueprint v2.2 闭环研究结果整理, 会话 34)}
\date{2026-08-09}

\begin{document}
\maketitle

\begin{abstract}
对 Dirichlet 加权弦 $-y''=\lambda\rho y$, 密度在单个区间 $(a,b)$ 上等于
$R>1$、其余处等于 $1$。本文完整证明 O3a: 对每个 $R>1$, 参数三角形
$\{(a,b):0<a<b<1\}$ 中恰有一个 sign-consistent good root, 且它必为对称根
$a+b=1$。证明机制是\emph{相位比刚性}: 高密度区间的传输能量守恒把两个界面
残差的共同零点转化为同一严格单调函数 $r_\tau(x)=J_m(\tau x)/J_m(x)$ 在两个
外侧相位上的等值, 从而所有 good root 必在反射固定线 $a+b=1$ 上; 对称线上
精确降维为单变量 $\widetilde F(c)$ 的零点问题, 其导数符号 $\widetilde F'(c)<0$
于 $0<c<1/2$ 由完全初等解析估计证明: 含 $\partial_qM_2<0$ 在 $D$ 上的完全解析
证明、C4 区间段 $K>0$ 的完全解析证明、$J_1^{(2)}$ 在真曲线区域上的完全解析
下界 (定理 \ref{thm:j1e1}) 与 $J_2^{(2)}$ 在盒 $[0.655,1.0472]\times[1,2]$
上的完全解析上界 (W-分解链, 定理 \ref{thm:j2e1}), 不再夹取隐式相位根; 其余
参数区域严格为负, 故对称线上恰有一个 good root。伴随命题 $(LOG)$
( $\frac{d}{dc}\log(\widetilde M_{f1}/\widetilde M_{f2})<0$) 由定理
\ref{thm:LOG} 纯解析证明。全文只使用\emph{严格解析证明} (E1), 与\emph{普通数值扫描}
(E3, 仅交叉检验) 明确区分并逐一标注; 早期版本中的四族叶盒证书 (J1 16 叶盒、J2 67 叶盒、
C4 200 叶盒与 $(LOG)$ 128 叶盒) 与单变量事实验证器均已解析化移除: $J_2^{(2)}$ 证明所需的
55 项单变量事实由注 \ref{rem:env} 的有理包络方法逐项给出 E1 证明, 证书总表见附录
\ref{app:cert}; 历史与复现合同见第 \ref{sec:certs} 节。
\end{abstract}

\tableofcontents

\section{问题、定义与主结论}

\subsection{加权弦势垒族}

固定 $R>1$ 与 $0<a<b<1$, 定义分段常数密度
\begin{equation}
	\rho_{a,b}(x)=
	\begin{cases}
		R, & a<x<b,\\
		1, & 0<x<a\ \text{或}\ b<x<1,
	\end{cases}
\end{equation}
考虑正则 Dirichlet Sturm--Liouville 问题
\begin{equation}
	-y''=\lambda\rho_{a,b}\,y,\qquad y(0)=y(1)=0.
	\label{eq:sl}
\end{equation}
记其前两个本征值为 $0<\lambda_1<\lambda_2$, 并用斜率归一化固定本征函数:
\begin{equation}
	y_k(0)=0,\qquad y_k'(0)=1,\qquad k=1,2.
\end{equation}
令
\begin{equation}
	n_k:=\int_0^1\rho_{a,b}(x)\,y_k(x)^2\,dx>0,\qquad
	\hat y_k:=\frac{y_k}{\sqrt{n_k}},
\end{equation}
定义
\begin{equation}
	f_{a,b}(x):=\lambda_1\hat y_1(x)^2-\lambda_2\hat y_2(x)^2,
	\label{eq:f}
\end{equation}
以及两个界面残差
\begin{equation}
	R_1(a,b):=f_{a,b}(a),\qquad R_2(a,b):=f_{a,b}(b).
\end{equation}
由于第一本征函数在开区间内严格为正, 还可定义
\begin{equation}
	v_{a,b}(x):=\frac{y_2(x)}{y_1(x)},\qquad 0<x<1.
\end{equation}

\begin{definition}[sign-consistent good root]\label{def:goodroot}
若 $(a,b)$ 满足
\begin{equation}
	R_1(a,b)=R_2(a,b)=0,\qquad v_{a,b}(a)>0,\qquad v_{a,b}(b)<0,
\end{equation}
则称其为一个 sign-consistent good root。
\end{definition}

\begin{theorem}[O3a 与更强的 exact-one 结论]\label{thm:main}
对每个有限实数 $R>1$, 参数三角形
\begin{equation}
	\{(a,b):0<a<b<1\}
\end{equation}
中恰有一个 sign-consistent good root。该根必为对称根:
\begin{equation}
	a+b=1.
\end{equation}
特别地, 任意两个 sign-consistent good roots 必相同, 这就是 canonical O3a 的
at-most-one 结论。
\end{theorem}

\begin{remark}[与分支函数 $g_1,g_2,h$ 的关系]\label{rem:branch}
证明直接作用于定义 \ref{def:goodroot} 的全部 good-root 集合, 不假设 $R_1=0$
或 $R_2=0$ 的某个分支是单图, 也不使用 $h=g_1-g_2$ 的全局导数符号。任何代表
genuine good root 的分支交点都必须是本文得到的唯一对称点; 只满足一个残差或
不满足符号条件的额外 sheet 不属于本定理的对象。
\end{remark}

\subsection{证明结构}

任意 good root $\Rightarrow$ 传输能量与残差消元 $\Rightarrow$ $a+b=1$;
对称线精确降维 $\Rightarrow$ $\widetilde F$ 恰有一零点 $\Rightarrow$ 全域恰一
good root。

\begin{remark}[证据标注约定]\label{rem:evidence}
全文证据分两类, 按下述约定标注与使用。
\begin{enumerate}
	\item[(E1)] \emph{严格解析证明}: 由闭式恒等式、初等不等式、单调性与精确
	有理常数界组成, 构成定理成立性的依据。单变量事实的 E1 证明由第
	\ref{ss:j2e1} 子节注 \ref{rem:env} 的\emph{有理包络方法}给出, 证书总表见
	附录 \ref{app:cert} (表 \ref{tab:envprims}--\ref{tab:envderiv});
	\item[(E3)] \emph{普通数值扫描}: 高精度浮点采样, 只用于交叉检验与探索,
	不作为任何定理的结论依据。
\end{enumerate}
结论陈述本身只依赖 (E1)。
\end{remark}

\section{所需的一维谱事实}

\begin{lemma}[前两模态的符号与零点]\label{lem:modes}
在斜率归一化下:
\begin{enumerate}
	\item[(i)] $y_1(x)>0$ 于 $0<x<1$, 且 $y_1'(1)<0$;
	\item[(ii)] $y_2$ 在 $(0,1)$ 中恰有一个简单零点 $z$, 并且
	\begin{equation}
		y_2>0\ \text{于}\ (0,z),\qquad y_2<0\ \text{于}\ (z,1),\qquad y_2'(1)>0.
	\end{equation}
\end{enumerate}
\end{lemma}

\begin{proof}
正则分离边界 Sturm--Liouville 问题具有离散、严格递增的简单本征值。这里列出
本文实际使用部分的直接理由。

第一 Rayleigh 商极小元可替换为其绝对值而不增加能量; 若它在内部为零, 则一维强
最大值原理和常微分方程唯一性迫使其恒为零。因此第一本征函数可取严格正。斜率
归一化 $y_1'(0)=1$ 正好选择这个正号。

第二本征函数与正的第一本征函数关于权 $\rho_{a,b}$ 正交, 故必变号。若 $y_2$
至少有两个内部零点, 则至少有三个相邻结区间。把 $y_2$ 限制在每个结区间并在外部
延为零, 可得到三个支集互不相交、Rayleigh 商都等于 $\lambda_2$ 的函数。它们张成
的三维空间经 min--max 原理给出 $\lambda_3\le\lambda_2$, 与简单性矛盾。故内部零点
恰有一个。

任一本征函数的零点都是简单的, 因为 $y(x_0)=y'(x_0)=0$ 会由初值问题唯一性推出
$y\equiv0$。同一本征值的两个 Dirichlet 解的 Wronskian 为常数, 且在端点为零, 故
二者成比例; 这也给出简单性。最后, 端点导数的符号由各模态靠近 $x=1$ 时的固定
符号立即得到。
\end{proof}

\begin{remark}
后文只使用引理 \ref{lem:modes}、本征值次序 $s_1<s_2$ 以及分段常系数方程的显式
解; 不需要任何关于残差曲线连通性的谱结论。
\end{remark}

\section{任意 good root 必为对称根}

\subsection{相位变量与精确相位范围}

固定 $R>1$ 和任意 good root $(a,b)$。置
\begin{equation}
	m:=\sqrt R>1,\qquad s_k:=\sqrt{\lambda_k},\qquad
	\tau:=\frac{s_2}{s_1}>1,
	\qquad \alpha:=s_1a,\qquad \beta:=s_1(1-b).
\end{equation}

\begin{lemma}[外侧相位范围]\label{lem:phases}
在任意 good root 处,
\begin{equation}
	0<\tau\alpha=s_2a<\pi,\qquad 0<\tau\beta=s_2(1-b)<\pi.
\end{equation}
等价地, $\alpha,\beta\in(0,\pi/\tau)$。
\end{lemma}

\begin{proof}
由 $y_1>0$ 和 good-root 符号,
\begin{equation}
	y_2(a)>0,\qquad y_2(b)<0.
\end{equation}
引理 \ref{lem:modes} 的唯一零点 $z$ 因而满足 $a<z<b$。

左侧密度为 $1$, 斜率归一化给出
\begin{equation}
	y_2(x)=\frac{\sin(s_2x)}{s_2},\qquad 0\le x\le a.
\end{equation}
若 $s_2a\ge\pi$, 则 $\pi/s_2\le a$ 是一个位于 $z$ 之前的零点, 矛盾。因此
$0<s_2a<\pi$。

右侧同样有某个非零常数 $C_2=-y_2'(1)$ 使
\begin{equation}
	y_2(x)=C_2\frac{\sin(s_2(1-x))}{s_2},\qquad b\le x\le1.
\end{equation}
若 $s_2(1-b)\ge\pi$, 则 $1-\pi/s_2\in[b,1)$ 给出 $z$ 之后的另一个零点, 仍矛盾。
故第二个不等式也成立。
\end{proof}

\subsection{高密度区间的传输不变量}

对任意谱参数 $s>0$, 在中间区间 $(a,b)$ 上定义状态
\begin{equation}
	U(x):=\binom{sy(x)}{y'(x)}.
\end{equation}
令
\begin{equation}
	\theta:=ms(b-a).
\end{equation}
常系数方程 $-y''=m^2s^2y$ 给出精确传输矩阵
\begin{equation}
	U(b)=P_m(\theta)U(a),\qquad
	P_m(\theta)=
	\begin{pmatrix}
		\cos\theta & \dfrac{\sin\theta}{m}\\[2mm]
		-m\sin\theta & \cos\theta
	\end{pmatrix}.
	\label{eq:transfer}
\end{equation}

\begin{lemma}[传输能量守恒]\label{lem:energy}
二次型
\begin{equation}
	Q_m(X,Y):=m^2X^2+Y^2
\end{equation}
被 $P_m(\theta)$ 保持。因而若
\begin{equation}
	W_m(x):=\cos^2x+m^2\sin^2x,\qquad J_m(x):=\frac{\sin^2x}{W_m(x)},
\end{equation}
则每个本征模都满足
\begin{equation}
	\frac{y_s(b)^2}{y_s(a)^2}=\frac{J_m(s(1-b))}{J_m(sa)}.
	\label{eq:ratio}
\end{equation}
\end{lemma}

\begin{proof}
直接展开得到
\begin{equation}
	m^2\left(\cos\theta\,X+\frac{\sin\theta}{m}Y\right)^2
	+(-m\sin\theta\,X+\cos\theta\,Y)^2=m^2X^2+Y^2.
\end{equation}

左界面上
\begin{equation}
	U(a)=\binom{\sin(sa)}{\cos(sa)}.
\end{equation}
右侧 Dirichlet 解写成
\begin{equation}
	y(x)=C_s\frac{\sin(s(1-x))}{s},
\end{equation}
故
\begin{equation}
	U(b)=C_s\binom{\sin(s(1-b))}{-\cos(s(1-b))}.
\end{equation}
守恒给出
\begin{equation}
	W_m(sa)=C_s^2\,W_m(s(1-b)).
\end{equation}
再用
\begin{equation}
	y(a)=\frac{\sin(sa)}{s},\qquad y(b)=C_s\frac{\sin(s(1-b))}{s}
\end{equation}
即得 \eqref{eq:ratio}。注意 $W_m>0$, 且此处没有使用任何本征函数积分归一化或
secular equation 的导数。
\end{proof}

\subsection{残差消元}

在 good root 处, 由 $R_1=R_2=0$, 有
\begin{equation}
	\frac{\lambda_1y_1(a)^2}{n_1}=\frac{\lambda_2y_2(a)^2}{n_2}>0,\qquad
	\frac{\lambda_1y_1(b)^2}{n_1}=\frac{\lambda_2y_2(b)^2}{n_2}>0.
\end{equation}
相位范围保证所除界面值非零。把第二式除以第一式, 得到
\begin{equation}
	\frac{y_1(b)^2}{y_1(a)^2}=\frac{y_2(b)^2}{y_2(a)^2}.
	\label{eq:elim}
\end{equation}
因此全部 $\lambda_k,n_k$ 因子都在同一模态内精确消去。对 \eqref{eq:elim} 的两侧
应用 \eqref{eq:ratio}, 得到
\begin{equation}
	r_\tau(\alpha)=r_\tau(\beta),\qquad
	r_\tau(x):=\frac{J_m(\tau x)}{J_m(x)},\qquad 0<x<\frac{\pi}{\tau}.
	\label{eq:rtaueq}
\end{equation}

\subsection{相位比的严格单调性}

\begin{lemma}[相位比单射]\label{lem:rtau}
对任意 $m\ge1$ 与 $\tau>1$, 函数 $r_\tau$ 在 $(0,\pi/\tau)$ 上严格递减。
\end{lemma}

\begin{proof}
写
\begin{equation}
	q_0:=m^2-1\ge0,\qquad W(x)=1+q_0\sin^2x,\qquad J(x)=\frac{\sin^2x}{W(x)}.
\end{equation}
对 $0<x<\pi$,
\begin{equation}
	\frac{d}{dx}\log J(x)=\frac{2\cot x}{W(x)}.
\end{equation}
定义
\begin{equation}
	\Psi(x):=\frac{x\cot x}{W(x)}.
\end{equation}
将 $\Psi'$ 乘以正量 $W(x)^2\sin^2x$, 直接化简为
\begin{equation}
	W(x)^2\sin^2x\,\Psi'(x)
	=\sin x\cos x-x
	+q_0\sin^2x\bigl(\sin x\cos x-x(1+2\cos^2x)\bigr).
	\label{eq:psi}
\end{equation}
令
\begin{equation}
	G(x):=x-\sin x\cos x.
\end{equation}
则 $G(0)=0$, 且
\begin{equation}
	G'(x)=2\sin^2x>0,\qquad 0<x<\pi.
\end{equation}
所以 $\sin x\cos x-x<0$。第二个括号等于
\begin{equation}
	-G(x)-2x\cos^2x<0.
\end{equation}
故 $\Psi'(x)<0$ 于整个 $(0,\pi)$。

若 $0<x<\pi/\tau$, 则 $x,\tau x\in(0,\pi)$, 从而
\begin{equation}
	\frac{d}{dx}\log r_\tau(x)
	=2\tau\frac{\cot(\tau x)}{W(\tau x)}-2\frac{\cot x}{W(x)}
	=\frac{2}{x}\bigl(\Psi(\tau x)-\Psi(x)\bigr)<0.
\end{equation}
由于 $r_\tau>0$, 它本身严格递减。
\end{proof}

\begin{theorem}[全 residual-sheet 的对称刚性]\label{thm:rigidity}
任意 sign-consistent good root 都满足 $a+b=1$。
\end{theorem}

\begin{proof}
由引理 \ref{lem:phases}, $\alpha,\beta$ 都位于引理 \ref{lem:rtau} 的单射区间。
等式 \eqref{eq:rtaueq} 因而推出 $\alpha=\beta$。又 $s_1>0$, 所以
\begin{equation}
	s_1a=s_1(1-b)\quad\Longrightarrow\quad a=1-b.
\end{equation}
起点是任意 good root, 因此证明覆盖每一条 residual sheet, 无需图结构、连通性、
Jacobian 非退化性或延拓假设。
\end{proof}

\section{对称线上的精确一维降维}

\subsection{对称参数与奇偶半问题}

现在置
\begin{equation}
	a=\xi,\qquad b=1-\xi,\qquad 0<\xi<\frac12,
\end{equation}
并改记
\begin{equation}
	q:=\sqrt R>1,\qquad \eta:=\frac12-\xi,\qquad c:=\frac{q\eta}{\xi}.
\end{equation}
于是
\begin{equation}
	\xi=\frac{q}{2(c+q)},\qquad \frac{dc}{d\xi}=-\frac{q}{2\xi^2}.
	\label{eq:xi}
\end{equation}
这给出严格递减双射
\begin{equation}
	\xi\in(0,1/2)\quad\longleftrightarrow\quad c\in(\infty,0).
\end{equation}

密度关于 $x=1/2$ 对称。由本征值简单性, 第一模态关于中点为偶函数, 第二模态为
奇函数。故在半区间 $[0,1/2]$ 上, 第一模态在中点满足 Neumann 条件, 第二模态满足
Dirichlet 条件。

令
\begin{equation}
	\alpha_k:=s_k\xi,\qquad qs_k\eta=c\alpha_k,
\end{equation}
以及
\begin{equation}
	\Phi_q(x):=\cos^2x+q^2\sin^2x.
\end{equation}
界面匹配给出
\begin{equation}
	\tan\alpha_1\tan(c\alpha_1)=\frac1q,\qquad
	q\tan\alpha_2+\tan(c\alpha_2)=0.
	\label{eq:match}
\end{equation}
为固定正确相位支, 定义
\begin{equation}
	E(x)=\arctan\frac{1}{q\tan x},\qquad 0<x<\frac{\pi}{2},
\end{equation}
以及连续分段函数
\begin{equation}
	O(x)=
	\begin{cases}
		\pi-\arctan(q\tan x), & 0<x<\pi/2,\\
		\pi/2, & x=\pi/2,\\
		\arctan(-q\tan x), & \pi/2<x<\pi.
	\end{cases}
\end{equation}

\begin{lemma}[真实相位落在主支]\label{lem:phasebranch}
对对称线上的前两个相位成立
\begin{equation}
	\alpha_1\in(0,\pi/2),\qquad \alpha_2\in(0,\pi),
	\label{eq:alpharange}
\end{equation}
且
\begin{equation}
	E(\alpha_1)=c\alpha_1,\qquad O(\alpha_2)=c\alpha_2.
	\label{eq:phasebranch}
\end{equation}
\end{lemma}

\begin{proof}
对问题 $-y''=\lambda\rho y$ 取 $s:=\sqrt\lambda$, 定义 Prufer 相位 $\theta$ 为
$y=r\sin\theta$, $y'=sr\cos\theta$, $\theta(0)=0$。则
\begin{equation}
	\theta'=s(\cos^2\theta+\rho\sin^2\theta)>0,
	\label{eq:prufer}
\end{equation}
故 $\theta$ 严格递增; 在左侧区域 $\rho=1$ 上 $y_k=\sin(s_kx)/s_k$, 故
$\theta_k(x)=s_kx$, 特别 $\alpha_k=\theta_k(\xi)$。

第一模态关于中点偶对称, 第二模态奇对称 ($y_1'(1/2)=0$, $y_2(1/2)=0$)。由引理
\ref{lem:modes}, $y_1>0$ 于 $(0,1)$, 而 $y_2$ 的唯一零点即 $z=1/2$, 故
$y_2>0$ 于 $(0,1/2)$。于是 $\sin\theta_k>0$ 于 $(0,1/2)$ ($k=1,2$); 由
$\theta_k(0)=0$ 与严格递增得 $\theta_1(1/2)=\pi/2$, $\theta_2(1/2)=\pi$, 从而
\begin{equation}
	\alpha_1=\theta_1(\xi)\in(0,\pi/2),\qquad
	\alpha_2=\theta_2(\xi)\in(0,\pi).
\end{equation}

在中间半区间 $(\xi,1/2)$ 上, 中点条件给出
\begin{equation}
	y_1(x)=A_1\cos(ms_1(x-1/2)),\qquad
	y_2(x)=A_2\sin(ms_2(x-1/2)),
\end{equation}
其幅角 $ms_k(x-1/2)$ 从 $x=\xi$ 处的 $-c\alpha_k$ 递增到 $x=1/2$ 处的 $0$。由
$y_1,y_2>0$ 于 $(\xi,1/2)$ 及幅角以 $0$ 为右端点, 得
\begin{equation}
	c\alpha_1\in(0,\pi/2),\qquad c\alpha_2\in(0,\pi).
\end{equation}

界面匹配 \eqref{eq:match} 的偶支给出 $\tan(c\alpha_1)=1/(q\tan\alpha_1)$, 与
$c\alpha_1\in(0,\pi/2)$ 合起来即 $c\alpha_1=E(\alpha_1)$。奇支在
$\alpha_2\ne\pi/2$ 时给出 $\tan(c\alpha_2)=-q\tan\alpha_2$: 若
$\alpha_2<\pi/2$, 则 $\tan(c\alpha_2)<0$, 结合 $c\alpha_2\in(0,\pi)$ 得
$c\alpha_2\in(\pi/2,\pi)$, 而 $O(\alpha_2)=\pi-\arctan(q\tan\alpha_2)$ 属同一区间
且正切相同, 故相等; $\alpha_2>\pi/2$ 时对称地得
$c\alpha_2=O(\alpha_2)=\arctan(-q\tan\alpha_2)\in(0,\pi/2)$。若
$\alpha_2=\pi/2$, 则 $y_2'(\xi)=\cos\alpha_2=0=A_2ms_2\cos(c\alpha_2)$ 给出
$\cos(c\alpha_2)=0$, 由 $c\alpha_2\in(0,\pi)$ 得 $c\alpha_2=\pi/2=O(\pi/2)$。
\end{proof}

于是正确的前两相位由
\begin{equation}
	E(\alpha_1)=c\alpha_1,\quad \alpha_1\in(0,\pi/2),\qquad
	O(\alpha_2)=c\alpha_2,\quad \alpha_2\in(0,\pi)
	\label{eq:phaseeq}
\end{equation}
唯一确定: 事实上
\begin{equation}
	E'(x)=O'(x)=-\frac{q}{\Phi_q(x)}<0,
\end{equation}
而 $cx$ 严格递增, 故 $E-cx$ 与 $O-cx$ 严格递减, 两条曲线与直线至多交一次; 引理
\ref{lem:phasebranch} 表明真实相位正是其交点。隐式求导得到
\begin{equation}
	\alpha_k'(c)=-\frac{\alpha_k\Phi_q(\alpha_k)}{q+c\Phi_q(\alpha_k)}.
	\label{eq:alphap}
\end{equation}

\subsection{谱隙导数与 canonical residual}

由 \eqref{eq:xi},
\begin{equation}
	s_k=\frac{\alpha_k}{\xi}=\frac{2(c+q)}{q}\alpha_k,
\end{equation}
故谱隙
\begin{equation}
	D(c):=\lambda_2-\lambda_1=\frac{4(c+q)^2}{q^2}\bigl(\alpha_2(c)^2-\alpha_1(c)^2\bigr).
	\label{eq:D}
\end{equation}
定义
\begin{equation}
	\widetilde M_f(x;c):=\frac{x^2\sin^2x}{q+c\Phi_q(x)},\qquad
	\widetilde F_e(c):=\widetilde M_f(\alpha_1(c);c)-\widetilde M_f(\alpha_2(c);c).
	\label{eq:Fe}
\end{equation}

\begin{lemma}[精确正因子降维]\label{lem:dimred}
若
\begin{equation}
	S_R(\xi):=R_1(\xi,1-\xi)=R_2(\xi,1-\xi),
\end{equation}
则
\begin{equation}
	S_R(\xi)=\frac{2(c+q)}{\xi^2}\widetilde F_e(c).
	\label{eq:dimred}
\end{equation}
因此 $S_R$ 与 $\widetilde F_e$ 具有完全相同的符号和零点。
\end{lemma}

\begin{proof}
将 \eqref{eq:alphap} 代入 \eqref{eq:D}。对任一相位 $x=\alpha_k$,
\begin{equation}
	\begin{aligned}
		\frac{d}{dc}\bigl((c+q)^2x^2\bigr)
		&=2(c+q)x^2+2(c+q)^2xx'\\
		&=-\frac{2q(c+q)(q^2-1)x^2\sin^2x}{q+c\Phi_q(x)}.
	\end{aligned}
\end{equation}
两模相减后得到
\begin{equation}
	D_c=\frac{8(c+q)(q^2-1)}{q}\widetilde F_e(c).
	\label{eq:Dc}
\end{equation}

另一方面, 令 $\hat y_k$ 按加权 $L^2$ 范数归一化。移动两个对称界面的
Feynman--Hellmann 公式为
\begin{equation}
	\frac{d\lambda_k}{d\xi}=(R-1)\lambda_k\bigl(\hat y_k(\xi)^2+\hat y_k(1-\xi)^2\bigr).
\end{equation}
奇偶性使平方值相等, 故
\begin{equation}
	D_\xi=2(q^2-1)\bigl(\lambda_2\hat y_2(\xi)^2-\lambda_1\hat y_1(\xi)^2\bigr)
	=-2(q^2-1)S_R(\xi).
\end{equation}
再由 $D_\xi=D_c\,c_\xi$、\eqref{eq:xi} 与 \eqref{eq:Dc} 即得 \eqref{eq:dimred}。
其前因子严格为正。
\end{proof}

\section{对称标量残差的唯一零点}

\subsection{\texorpdfstring{容易处理的区域 $c\ge1/2$}{容易处理的区域 c ge 1/2}}

固定 $c>0$, 记
\begin{equation}
	\varphi_c(x):=\frac{x^2\sin^2x}{q+c\Phi_q(x)}.
\end{equation}
在 $0<x<\pi/2$ 上,
\begin{equation}
	\frac{d}{dx}\log\varphi_c(x)
	=\frac2x+2\cot x-\frac{2c(q^2-1)\sin x\cos x}{q+c\Phi_q(x)}.
\end{equation}
由于
\begin{equation}
	\frac{c}{q+c\Phi_q(x)}<\frac{1}{\Phi_q(x)}
\end{equation}
以及
\begin{equation}
	\frac{(q^2-1)\sin x\cos x}{\Phi_q(x)}
	=\frac{(q^2-1)\tan x}{1+q^2\tan^2x}<\cot x,
\end{equation}
故 $\varphi_c$ 严格递增。

若 $c\ge1$, 则 $\alpha_1<\alpha_2\le\pi/2$, 从而
\begin{equation}
	\widetilde F_e(c)<0.
\end{equation}
若 $1/2\le c<1$, 写
\begin{equation}
	\alpha_2=\pi-\gamma,\qquad 0<\gamma<\pi/2.
\end{equation}
奇相位方程等价于
\begin{equation}
	E(\gamma)-c\gamma=\pi(1/2-c)\le0.
\end{equation}
函数 $E(x)-cx$ 严格递减, 并在 $x=\alpha_1$ 为零, 因此 $\gamma\ge\alpha_1$。于是
\begin{equation}
	\widetilde M_f(\alpha_2;c)=\left(\frac{\pi-\gamma}{\gamma}\right)^2\varphi_c(\gamma)
	>\varphi_c(\alpha_1)=\widetilde M_f(\alpha_1;c),
\end{equation}
而端点 $c=1$ 也由连续性和 $\alpha_1<\alpha_2=\pi/2$ 得到严格不等式。所以
\begin{equation}
	\widetilde F_e(c)<0,\qquad c\ge\frac12.
	\label{eq:Feneg}
\end{equation}
另一方面, 由相位曲线直接看出
\begin{equation}
	\alpha_1(c)\to\frac{\pi}{2},\qquad \alpha_2(c)\to\pi\qquad(c\to0^+),
\end{equation}
因而
\begin{equation}
	\widetilde F_e(c)\to\frac{\pi^2}{4q}>0.
	\label{eq:Fepos}
\end{equation}

\subsection{关键引理}

\begin{theorem}[KEY LEMMA, 全参数不等式]\label{thm:keylemma}
对每个 $q>1$ 与 $0<c<1/2$,
\begin{equation}
	\widetilde F_e'(c)<0.
	\label{eq:keylemma}
\end{equation}
\end{theorem}

本节给出该引理的完整证明链。伴随命题 $(LOG)$ 由定理 \ref{thm:LOG} 完全解析证明 (E1, 见第 \ref{ss:log} 子节); I3 的二维参数化不等式由定理 \ref{thm:j1e1} 与 \ref{thm:j2e1} 完全解析化;
半无限参数尾部与 C4 曲线段均由解析估计处理。

\subsubsection{沿相位曲线的对数导数}

对 $x=\alpha_k(c)$, 由 \eqref{eq:alphap} 直接计算
\begin{equation}
	\frac{d}{dc}\log\widetilde M_f(\alpha_k(c);c)=G(\alpha_k(c);c),
\end{equation}
其中
\begin{equation}
	G(x;c)=-\frac{\Phi_q(x)\bigl[3+2x\cot x\bigr]}{q+c\Phi_q(x)}
	+\frac{2cx\Phi_q(x)(q^2-1)\sin x\cos x}{\bigl[q+c\Phi_q(x)\bigr]^2}.
	\label{eq:G}
\end{equation}
记 $G_k:=G(\alpha_k(c);c)$, 则
\begin{equation}
	\widetilde F_e'=\widetilde M_{f1}G_1-\widetilde M_{f2}G_2.
	\label{eq:Fep}
\end{equation}

\begin{lemma}\label{lem:G1}
$G_1<0$ 对全部 $q>1$, $c>0$ 成立。
\end{lemma}

\begin{proof}
因 $0<\alpha_1<\pi/2$, 有
\begin{equation}
	(q^2-1)\sin x\cos x\le\Phi_q(x)\cot x.
\end{equation}
令 $D=q+c\Phi_q(x)>0$。将 \eqref{eq:G} 乘以 $D^2/\Phi_q(x)>0$, 并用上式估计正项,
得到
\begin{equation}
	\frac{D^2}{\Phi_q}G\le-\bigl[3+2x\cot x\bigr]D+2cx\Phi_q\cot x
	=-3D-2qx\cot x<0.
\end{equation}
\end{proof}

若 $G_2\ge0$, 由 \eqref{eq:Fep} 和 $\widetilde M_{f1},\widetilde M_{f2}>0$ 立即得到
$\widetilde F_e'<0$。因此只需把
\begin{equation}
	B:=\{(q,c):G_2<0,\ q>1,\ 0<c<1/2\}
\end{equation}
限制到一个紧盒, 并在该盒上证明 $\widetilde F_e'<0$。

\subsubsection{\texorpdfstring{把 $G_2<0$ 限制到紧盒}{把 G2<0 限制到紧盒}}

写
\begin{equation}
	\gamma:=\pi-\alpha_2\in(0,\pi/2),\qquad w:=q\tan\gamma=\tan(c\alpha_2),
\end{equation}
并令
\begin{equation}
	A:=\alpha_2=\pi-\arctan(w/q),\qquad c=\frac{\arctan w}{A}.
\end{equation}
在 $0<c<1/2$ 上
\begin{equation}
	0<w<\sqrt{2q+1},
\end{equation}
且 $w$ 随 $c$ 严格增加。定义
\begin{equation}
	\mathrm{IN}(q,w):=(q^2+w^2)A\bigl[2Aq-3w+2\arctan w\bigr]
	-3wq(1+w^2)\arctan w.
	\label{eq:IN}
\end{equation}
精确代数恒等式给出
\begin{equation}
	\mathrm{IN}=G_2\cdot\mathrm{POS},
\end{equation}
其中
\begin{equation}
	\mathrm{POS}=\frac{\bigl[q+c\Phi_q(\alpha_2)\bigr]^2A(q^2+w^2)w}{\Phi_q(\alpha_2)\,q};
\end{equation}
故 $G_2$ 与 $\mathrm{IN}$ 同号。

对 $w$ 求导:
\begin{equation}
	M_2(q,w):=\partial_w\mathrm{IN}
	=4A^2wq-7Aq^2-9Aw^2+\frac{2A(q^2+w^2)}{1+w^2}
	+\arctan w\bigl[4Aw-5q-9qw^2\bigr].
	\label{eq:M2}
\end{equation}

\begin{lemma}[$q=1$ 基线上的 $\partial_qM_2$]\label{lem:B1}
置 $g(w):=\partial_qM_2(1,w)$。则 $g(w)<0$ 对一切 $w\in[0,\sqrt3]$ 成立。
\end{lemma}

\begin{proof}
记 $t:=\arctan w$。直接化简给出
\begin{equation}
	g''(w)=-\frac{2\bigl(9w^6t+9w^5+27w^4t+24w^3+19w^2t+20\pi w^2+31w+t+4\pi\bigr)}{(1+w^2)^3}<0,
	\label{eq:gpp}
\end{equation}
故 $g'$ 在 $[0,\infty)$ 严格递减。又有
\begin{equation}
	\begin{split}
		g'(0)&=4\pi^2>0,\\
		g'(\sqrt3)&=\frac{16\pi^2}{9}-\frac{41\sqrt3\,\pi}{6}-15
		\le\frac{16(22/7)^2}{9}-\frac{41(17/10)(157/50)}{6}-15
		=-\frac{14957063}{441000}<0.
	\end{split}
\end{equation}
其中末步用 $\pi\le22/7$ 与 $\sqrt3\ge17/10$。于是 $g'$ 在 $(0,\sqrt3)$ 有唯一零点
$w_0$, 且 $g$ 恰在 $w_0$ 取得唯一最大值。

置 $b:=\arctan(4/5)$。$\arctan$ 的 Leibniz 交错级数在 $x=4/5$ 的部分和
$S_k:=\sum_{j=0}^{k}(-1)^jx^{2j+1}/(2j+1)$ 满足 $S_5<b<S_6$, 其中
\begin{equation}
	S_5=\frac{22739538548}{33837890625},\qquad
	S_6=\frac{7436856470852}{10997314453125},
\end{equation}
且直接有理比较给出
\begin{equation}
	\frac{67}{100}<S_5<b<S_6<\frac{17}{25}.
\end{equation}
精确化简给出
\begin{equation}
	g(4/5)=-\frac{19b}{25}-\frac{346\pi}{41}+\frac{16(b-\pi)^2}{5}-\frac{1076}{205},
	\qquad
	g'(4/5)=-\frac{2152b}{205}-\frac{2560\pi}{1681}+4(b-\pi)^2-\frac{15008}{1681}.
\end{equation}
对 $b\in(67/100,17/25)$ 与 $\pi\in(157/50,22/7)$ 有
\begin{equation}
	\partial_bg(4/5)<0,\quad \partial_\pi g(4/5)>0,\quad
	\partial_bg'(4/5)<0,\quad \partial_\pi g'(4/5)>0,
\end{equation}
故
\begin{equation}
	g(4/5)\le g(4/5;b=67/100,\pi=22/7)<0,
	\qquad
	g'(4/5)\ge g'(4/5;b=17/25,\pi=157/50)>3.3581>0.
\end{equation}
因此 $w_0\in(4/5,\sqrt3)$。由 $g''<0$ (严格凹), 图像位于切线之下, 故
\begin{equation}
	\begin{split}
		g(w_0)&\le g(4/5)+g'(4/5)(w_0-4/5)
		\le g(4/5)+g'(4/5)(\sqrt3-4/5)\\
		&\le g(4/5;b=67/100,\pi=22/7)+g'(4/5;b=67/100,\pi=22/7)\Bigl(\frac74-\frac45\Bigr)
		=-\frac{1054523}{114800}<0.
	\end{split}
\end{equation}
于是 $\max_{[0,\sqrt3]}g=g(w_0)<0$。
\end{proof}

\begin{lemma}[边界曲线 $w=\sqrt{2q+1}$]\label{lem:boundary}
对一切 $q\ge1$, 记 $w_b(q):=\sqrt{2q+1}$。则
\begin{equation}
	\partial_qM_2(q,w_b(q))<0,\qquad M_2(q,w_b(q))<0.
	\label{eq:boundary}
\end{equation}
\end{lemma}

\begin{proof}
置 $s:=w_b(q)=\sqrt{2q+1}\ge\sqrt3$ 与 $\theta:=\arctan(1/s)\in(0,\pi/6]$, 则
\begin{equation}
	q=\frac{\cos2\theta}{2\sin^2\theta},\qquad w=s=\cot\theta,\qquad
	\arctan\frac{w}{q}=2\theta,\qquad
	\arctan w=\frac{\pi}{2}-\theta,\qquad A=\pi-2\theta.
\end{equation}
把 \eqref{eq:M2} 代入后精确化简 (符号计算 diff 为零) 给出
\begin{equation}
	M_2(q,w_b(q))=2(2\theta-\pi)\cot^2\theta
	\Bigl[(2\theta-\pi)\cot\theta+\frac{2}{\sin^2\theta}\Bigr].
	\label{eq:Gtheta}
\end{equation}
因 $2\theta-\pi<0$ 且
\begin{equation}
	\Bigl[(2\theta-\pi)\cot\theta+\frac{2}{\sin^2\theta}\Bigr]\sin^2\theta
	=2-\Bigl(\frac{\pi}{2}-\theta\Bigr)\sin2\theta\ge2-\frac{\pi}{2}>0,
\end{equation}
得到 $M_2(q,w_b(q))<0$。

对 $\partial_qM_2$ 令 $z:=\tan\theta\in(0,1/\sqrt3]$ 与 $\beta:=\arctan z$。代入化简得
\begin{equation}
	\partial_qM_2(q,w_b(q))=\frac{N(z)}{2z^2(z^2+1)^2},
	\qquad
	N(z)=\beta^2P(z)+\beta Q(z)+R(z),
	\label{eq:Fz}
\end{equation}
其中
\begin{equation}
	P(z)=32z(z^2+1)^2,
\end{equation}
\begin{equation}
	Q(z)=-10z^6-32\pi z^5+42z^4-64\pi z^3+2z^2-32\pi z+46,
\end{equation}
\begin{equation}
	R(z)=5\pi z^6-10z^5+8\pi^2z^5-21\pi z^4-40z^3+16\pi^2z^3-\pi z^2-14z+8\pi^2z-23\pi.
\end{equation}
因 $P(z)>0$, $N$ 作为 $\beta\in[0,\pi/6]$ 的二次函数严格凸, 故其最大值在端点
$\beta=0$ 或 $\beta=\pi/6$ 取得:
\begin{equation}
	N(z)\le\max\bigl\{R(z),T(z)\bigr\},\qquad
	T(z):=\frac{\pi^2}{36}P(z)+\frac{\pi}{6}Q(z)+R(z).
\end{equation}
展开并整理
\begin{equation}
	T(z)=\frac{10\pi}{3}z^6+\Bigl(\frac{32\pi^2}{9}-10\Bigr)z^5-14\pi z^4
	+\Bigl(\frac{64\pi^2}{9}-40\Bigr)z^3-\frac{2\pi}{3}z^2
	+\Bigl(\frac{32\pi^2}{9}-14\Bigr)z-\frac{46\pi}{3},
\end{equation}
\begin{equation}
	R(z)=-23\pi+(8\pi^2-14)z+(16\pi^2-40)z^3+(8\pi^2-10)z^5+5\pi z^6
	-21\pi z^4-\pi z^2.
\end{equation}
在 $z\le1/\sqrt3\le10/17$ 与 $\pi\in(157/50,22/7)$ 下, 丢弃 $R,T$ 中所有负项并
对正系数用 $z\le10/17$、$\pi\le22/7$, 对负常数项用 $\pi\ge157/50$, 得到精确有理上界
\begin{equation}
	R(z)\le-\frac{262235520291}{59137044050}<0,
	\qquad
	T(z)\le-\frac{7282185739373}{266116698225}<0.
\end{equation}
故 $N(z)<0$, 即 $\partial_qM_2(q,w_b(q))<0$。
\end{proof}

\begin{lemma}[$M_2$]\label{lem:M2}
在
\begin{equation}
	D=\{(q,w):q>1,\ 0<w<\sqrt{2q+1}\}
\end{equation}
上, $M_2(q,w)<0$。因此 $\mathrm{IN}(q,w)$ 对 $w$ 严格递减。
\end{lemma}

\begin{proof}
证明分五部分。记 $A:=\pi-\arctan(w/q)$, $t:=\arctan w$。

\medskip
\noindent\textbf{(a) 基线 $q=1$.} 直接化简为
\begin{equation}
	M_2(1,w)=\pi h(w),\qquad h(w)=4w(\pi-\arctan w)-5-9w^2.
\end{equation}
有
\begin{equation}
	h''(w)=-\frac{8}{(1+w^2)^2}-18<0.
\end{equation}
用交错级数有理界检查 $h'(1/2)>0.1016$、$h'(0.53)<-0.52$, 故 $h'$ 仅有一个零点
$w^*\in(0.5,0.53)$。在极大点利用 $h'(w^*)=0$ 得
\begin{equation}
	h(w)\le h(w^*)=\frac{4w^{*2}}{1+w^{*2}}+9w^{*2}-5
	<13(0.53)^2-5=-1.3483<0.
\end{equation}
故 $M_2(1,w)<0$ 对一切 $w>0$ 成立。

\medskip
\noindent\textbf{(b) $\partial^2_qM_2<0$.} 直接化简 (符号计算 diff 为零) 给出
\begin{equation}
	\partial^2_qM_2=\frac{2N_2}{(q^2+w^2)^2(1+w^2)},
\end{equation}
\begin{equation}
	N_2=-A\bigl[7q^4w^2+5q^4+14q^2w^4+10q^2w^2-w^6-3w^4\bigr]
	-7q^3w^3-5q^3w-qw^5-4qw^4t+qw^3-4qw^2t.
	\label{eq:N2}
\end{equation}

\emph{(b1) 盒 $[1,20]\times[0,\sqrt3]$.} 因 $q\ge1$ 与 $w\le\sqrt3$, 有
$\arctan(w/q)\le\arctan w\le\pi/3$, 故 $A\ge2\pi/3$。又由 $w^2\le3$ 得 $w^6\le3w^4$,
所以
\begin{equation}
	7q^4w^2+5q^4+14q^2w^4+10q^2w^2-w^6-3w^4
	\ge(5+7w^2)q^4\ge5q^4.
\end{equation}
在 \eqref{eq:N2} 中丢弃全部非正项, 得
\begin{equation}
	N_2\le-\frac{2\pi}{3}(5+7w^2)q^4+qw^3
	\le-\frac{10\pi}{3}q^4+3\sqrt3\,q
	=q\Bigl[3\sqrt3-\frac{10\pi}{3}q^3\Bigr]<0,
\end{equation}
因为 $q\ge1$ 而 $3\sqrt3\le3(7/4)<10(157/50)/3\le10\pi/3$。故
$\partial^2_qM_2<0$ 于 $[1,20]\times[0,\sqrt3]$。

\emph{(b2) $D\cap\{w\ge\sqrt3\}$.} 由 $w^2<2q+1$ 得 $w^6<(2q+1)w^4$, 故
\begin{equation}
	14q^2w^4+10q^2w^2-w^6-3w^4
	>\bigl(14q^2-2q-4\bigr)w^4\ge8w^4>0,
\end{equation}
于是 \eqref{eq:N2} 中方括号为正, 记之为 $B_0>0$。又 $w\ge\sqrt3$ 时
$1-w^2-4wt\le1-w^2<0$, 故
\begin{equation}
	N_2\le-AB_0+qw^3(1-w^2-4wt)-4qw^2t<-AB_0<0.
\end{equation}
即 $\partial^2_qM_2<0$ 于 $D\cap\{w\ge\sqrt3\}$ (无 $q$ 上界)。

\medskip
\noindent\textbf{(c) $\partial_qM_2<0$ 于 $D\cap\{q\le20\}$.}
由 (b1), 对固定 $w\le\sqrt3$, $\partial_qM_2$ 在 $q\in[1,20]$ 严格递减, 故
\begin{equation}
	\partial_qM_2(q,w)<\partial_qM_2(1,w)=g(w)<0\qquad(w\le\sqrt3),
\end{equation}
其中末步为引理 \ref{lem:B1}。对 $\sqrt3\le w\le\sqrt{41}$, 由 (b2),
$\partial_qM_2$ 在 $q\in[m(w),20]$ 严格递减, 其中 $m(w):=(w^2-1)/2\ge1$ 满足
$w=\sqrt{2m(w)+1}$。于是
\begin{equation}
	\partial_qM_2(q,w)<\partial_qM_2(m(w),w)<0,
\end{equation}
末步为引理 \ref{lem:boundary} 的 \eqref{eq:boundary} 第一式。

\medskip
\noindent\textbf{(d) 尾部 $q\ge20$.} 利用
\begin{equation}
	A\le\pi,\qquad A\ge\pi-\frac{\sqrt{2q+1}}{q},\qquad
	w\le\sqrt{2q+1},\qquad \arctan w\le\frac{\pi}{2}
\end{equation}
得到
\begin{equation}
	\partial_qM_2\le B(q),
\end{equation}
\begin{equation}
	B(q)=(4\pi^2+14)\sqrt{2q+1}+8\pi\frac{2q+1}{q}+1+2\pi\frac{2q+1}{q^2}-10\pi q.
\end{equation}
有 $B(20)<-232.723$, 且\footnote{此数值界由有理包络验证: $B(20)=(4\pi^2+14)\sqrt{41}+1-\frac{183395}{1000}\pi$ 在盒 $\pi\in(\frac{314159}{100000},\frac{314160}{100000})$、$\sqrt{41}\in(\frac{640312}{100000},\frac{640313}{100000})$ 上满足 $\partial_\pi B(20)=8\pi\sqrt{41}-\frac{183395}{1000}<0$ 与 $\partial_{\sqrt{41}}B(20)=4\pi^2+14>0$, 故上界取在 $(\pi_{\min},\sqrt{41}_{\max})$, 精确有理比较给出 $B(20)\le-\frac{58180766243071047}{250000000000000}<-232.723$.}
\begin{equation}
	B'(q)\le\frac{4\pi^2+14}{\sqrt{41}}-10\pi<0.
\end{equation}
故半无限尾部同样严格为负: $\partial_qM_2<0$ 于 $D\cap\{q\ge20\}$。

\medskip
\noindent\textbf{(e) $M_2<0$ 于 $D$.} 若 $0<w\le\sqrt3$, 沿 $q$ 从 $1$ 积分
$\partial_qM_2$ 并由 (a) 得 $M_2(q,w)<M_2(1,w)<0$。若 $\sqrt3<w\le\sqrt{41}$,
从 $D$ 的边界 $q=m(w)$ 积分, 结合 (c)(d) 与引理 \ref{lem:boundary} 的
\eqref{eq:boundary} 第二式得 $M_2(q,w)<M_2(m(w),w)<0$。若 $w>\sqrt{41}$, 则
$q>20$。置 $t=w/q\le\sqrt{41}/20<0.33$, 并用 $q^2t^2>41$, 可估计
\begin{equation}
	\frac{M_2}{q^2}\le4\pi^2(0.33)-7(\pi-0.33)+\frac{2\pi(1+0.33^2)}{42}<0.
\end{equation}

五部分覆盖整个 $D$。
\end{proof}

\begin{lemma}[CORNER 与 C4]\label{lem:corner}
\begin{equation}
	G_2(1/2;q)\ge0\quad(q\ge2),\qquad G_2(0.4;q)\ge0\quad(q>1).
\end{equation}
\end{lemma}

\begin{proof}
\medskip
\noindent\textbf{CORNER.} 在 $c=1/2$ 上,
\begin{equation}
	\alpha_1=x,\qquad \alpha_2=\pi-x,\qquad
	\cos x=\frac{q}{q+1},\qquad \sin x=\frac{\sqrt{2q+1}}{q+1},
\end{equation}
并有精确公式
\begin{equation}
	G_2(1/2;q)=\frac{2q(q+1)(\pi-x-3\sin x)}{(2q+1)^{3/2}}.
\end{equation}
当 $q\ge2$ 时 $x\le\arccos(2/3)$。函数 $\pi-x-3\sin x$ 在所需区间的最小值位于
$q=2$, 而
\begin{equation}
	G_2(1/2;2)=\frac{12\bigl[\pi-\arccos(2/3)-\sqrt5\bigr]}{5\sqrt5}>0.
\end{equation}
最后一个严格不等式可由 $\pi>3.14$, $\sqrt5<2.24$ 及余弦交错级数上界直接核验。

\medskip
\noindent\textbf{C4.} 在 $c=0.4$ 上令 $v=\arctan w$, 则
\begin{equation}
	v\in[2\pi/7,2\pi/5),\qquad A=\frac52v,\qquad
	q=\frac{\tan v}{\tan(\pi-\frac52v)},\qquad w=\tan v.
\end{equation}
恒等式 $\mathrm{IN}=AK(v)$, 其中
\begin{equation}
	K(v)=(q^2+w^2)(5vq-3w+2v)-\frac65wq(1+w^2).
\end{equation}
下面给出 $K>0$ 于整个 $[2\pi/7,2\pi/5)$ 的纯初等解析证明, 完全取代原
200 叶盒证书与尾段处理。

设 $\omega=\pi-\frac52v\in(0,2\pi/7]$, $T=\tan\omega$, 则 $q=w/T$ 且
\begin{equation}
	K=q^2L,\qquad
	L(v)=(1+T^2)\Bigl(w\Bigl(\frac{5v}{T}-3\Bigr)+2v\Bigr)
	-\frac65T(1+w^2).
	\label{eq:Ldef}
\end{equation}
因 $v\ge\omega\iff v\ge2\pi/7$, 有 $q=w/T\ge1$; 故只要证 $L>0$ 即得 $K>0$。
下面证明 $L$ 在 $[2\pi/7,2\pi/5)$ 严格递增且 $L(2\pi/7)>0$。

先记录常数引理, 全部可由初等方式验证。Machin 级数给出
\begin{equation}
	\pi\in\Bigl(\frac{31415}{10000},\frac{31416}{10000}\Bigr),
\end{equation}
故
\begin{equation}
	\frac{2\pi}{7}>\frac{8975}{10000},\qquad
	\frac{3\pi}{10}<\frac{9425}{10000},\qquad
	\frac{3\pi}{10}>\frac{9424}{10000},\qquad
	\frac{2\pi}{5}<\frac{12567}{10000}.
\end{equation}
由 $\tan^2(3\pi/10)=1+\frac{2\sqrt5}{5}$ 及 $\sqrt5\in(2.2360,2.2361)$ 得
\begin{equation}
	\frac{13763}{10000}<\tan\frac{3\pi}{10}<\frac{13765}{10000},
\end{equation}
由 $\tan^2(2\pi/5)=5+2\sqrt5$ 得 $\tan(2\pi/5)<3078/1000$。又
$\tan(2\pi/7)$ 是多项式
\begin{equation}
	P(t)=t^6-21t^4+35t^2-7
\end{equation}
在 $(1,2)$ 中的唯一零点: $P$ 在 $[1,2]$ 严格递减 ($P'(t)=2t(3t^4-42t^2+35)$,
其第二因子在 $t^2\in[1,4]$ 上 $\le-4$) 且 $P(1)=8>0>P(2)=-139$, 而
$\tan(2\pi/7)\in(\tan\frac\pi4,\tan\frac\pi3)\subset(1,2)$ 满足 $P=0$
(它是 $\tan(7\theta)$ 的分子)。精确有理求值
\begin{equation}
	P\Bigl(\frac{1253}{1000}\Bigr)>0>P\Bigl(\frac{1254}{1000}\Bigr)
	\qquad\text{给出}\qquad
	\tan\frac{2\pi}{7}\in\Bigl(\frac{1253}{1000},\frac{1254}{1000}\Bigr).
\end{equation}

\medskip
\noindent\textbf{单调性.} 用商数法则代入 $w'=1+w^2$,
$T'=-\frac52(1+T^2)$ 直接求导得
\begin{equation}
	L'(v)=\frac{N(v)}{10T^2},
	\label{eq:Lp}
\end{equation}
其中
\begin{equation}
	N=125wv+50T\bigl(v(1+w^2)+w\bigr)+20T^2+c_3T^3+(20-125wv)T^4+c_5T^5,
\end{equation}
\begin{equation}
	c_3=50w^2v-24w^3+176w-50v,\qquad c_5=150w-100v.
\end{equation}
分两个区域证明 $N>0$。

\medskip
\noindent\textbf{区域 I: $v\in[2\pi/7,3\pi/10]$.}
此时 $T\in[1,\tan(2\pi/7)]$, $w\in[\tan(2\pi/7),\tan(3\pi/10)]$。除
$(20-125wv)T^4$ 外各项非负 ($wv\ge\frac{1253}{1000}\frac{8975}{10000}>\frac{20}{125}$), 故
\begin{equation}
	\begin{aligned}
		N\ge{}&125\frac{1253}{1000}\frac{8975}{10000}
		+50\Bigl(\frac{8975}{10000}\Bigl(1+\frac{1253^2}{10^6}\Bigr)
		+\frac{1253}{1000}\Bigr)+20\\
		&+\Bigl(50\frac{8975}{10000}\Bigl(\frac{1253^2}{10^6}-1\Bigr)
		+2\frac{1253}{1000}\Bigl(88-12\frac{13765^2}{10^8}\Bigr)\Bigr)
		+\Bigl(150\frac{1253}{1000}-100\frac{9425}{10000}\Bigr)\\
		&-125\frac{13765}{10000}\frac{9425}{10000}\Bigl(\frac{1254}{1000}\Bigr)^4
		=\frac{88146367488708279}{400000000000000}>0.
	\end{aligned}
\end{equation}
此处 $c_3=50v(w^2-1)+2w(88-12w^2)$ 的下界用 $v\ge\frac{8975}{10000}$,
$w\ge\frac{1253}{1000}$ 及 $88-12w^2\ge88-12\frac{13765^2}{10^8}>0$;
$c_5=150w-100v$ 的下界用 $w\ge\frac{1253}{1000}$, $v\le\frac{9425}{10000}$;
负项上界用 $w\le\frac{13765}{10000}$, $v\le\frac{9425}{10000}$,
$T\le\frac{1254}{1000}$。

\medskip
\noindent\textbf{区域 II: $v\in[3\pi/10,2\pi/5)$.}
此时 $T\le1$, 改写
\begin{equation}
	N=125wv(1-T^4)+20T^4+50T\bigl(v(1+w^2)+w\bigr)+20T^2+c_3T^3+c_5T^5.
\end{equation}
前三项、$20T^2$ 与 $c_5T^5$ 均非负 ($c_5\ge150\frac{13763}{10000}
-100\frac{12567}{10000}>0$)。仅需 $c_3\ge0$。若 $w\le\frac{27}{10}$, 则
\begin{equation}
	c_3\ge50v(w^2-1)+2w\Bigl(88-12w^2\Bigr)\ge
	2w\Bigl(88-12\frac{27^2}{10^2}\Bigr)>0;
\end{equation}
若 $w\ge\frac{27}{10}$, 则 $w\le\frac{3078}{1000}$, $v\ge\frac{9424}{10000}$,
且 $176w-24w^3$ 在 $[\frac{27}{10},\frac{3078}{1000}]$ 递减 (导数
$176-72w^2<0$), 于是
\begin{equation}
	c_3\ge50\frac{9424}{10000}\Bigl(\frac{27^2}{10^2}-1\Bigr)
	+\Bigl(176\frac{3078}{1000}-24\frac{3078^3}{10^9}\Bigr)
	=\frac{2160051043}{15625000}>0.
\end{equation}
故区域 II 上 $N>0$ (因 $20T^4>0$)。

\medskip
由 \eqref{eq:Lp} 得 $L'>0$, 即 $L$ 在 $[2\pi/7,2\pi/5)$ 严格递增。左端点
$q=1$, 故 $L(2\pi/7)=K(2\pi/7)$, 且
\begin{equation}
	L(2\pi/7)=(1+w^2)\Bigl(2\pi-\frac{21}{5}w\Bigr),
	\qquad w=\tan\frac{2\pi}{7},
\end{equation}
\begin{equation}
	2\pi-\frac{21}{5}w\ge2\frac{31415}{10000}
	-\frac{21}{5}\frac{1254}{1000}>0.
\end{equation}
故 $L(2\pi/7)\ge\frac{13058215729}{5000000000}>0$, 从而
$K=q^2L\ge L>0$ 于整个 $[2\pi/7,2\pi/5)$, C4 得证。
\end{proof}

由引理 \ref{lem:M2}, $\mathrm{IN}$ 随 $w$ 严格下降, 而 $w$ 随 $c$ 严格上升。
CORNER 表明 $q\ge2$ 时整个 $0<c<1/2$ 上 $G_2\ge0$; C4 表明全部 $q>1$ 在 $c\le0.4$
时 $G_2\ge0$。因此
\begin{equation}
	B\subset(1,2)\times(0.4,0.5).
	\label{eq:B}
\end{equation}

\subsubsection{沿相位曲线的二阶导数与二维参数化}

为证明 $\widetilde F_e''>0$ 于 $Q:=[1,2]\times[0.4,0.5]$ 成立, 先把
$\widetilde F_e''$ 化为显式函数的符号问题。记
\begin{equation}
	G_x(x;c):=\partial_xG(x;c),\qquad
	G_c(x;c):=\partial_cG(x;c),
\end{equation}
并沿相位曲线定义
\begin{equation}
	J(x;c):=G(x;c)^2+G_x(x;c)\alpha_k'(c)+G_c(x;c)
	=G(x;c)^2-\frac{x\Phi_q(x)}{q+c\Phi_q(x)}G_x(x;c)+G_c(x;c),
	\label{eq:J}
\end{equation}
其中末式由 \eqref{eq:alphap} 得到。由 \eqref{eq:Fep} 与
$\frac{d}{dc}\log\widetilde M_f(\alpha_k(c);c)=G(\alpha_k(c);c)$,
\begin{equation}
	\widetilde F_e''(q,c)=\widetilde M_{f1}J(\alpha_1;c)-\widetilde M_{f2}J(\alpha_2;c).
	\label{eq:Fep2raw}
\end{equation}

相位方程 \eqref{eq:phaseeq} 沿真实曲线可显式反解。对第一相位, 由
$E(x)=\arctan(1/(q\tan x))$,
\begin{equation}
	c=c_1(x,q):=\frac1x\arctan\frac1{q\tan x},\qquad 0<x<\frac\pi2,
	\label{eq:c1}
\end{equation}
即对每个 $(q,c)\in Q$, 真实点 $x=\alpha_1(q,c)$ 满足 $c=c_1(x,q)$; 定义
\begin{equation}
	J_1^{(2)}(x,q):=J(x;c_1(x,q)).
\end{equation}
对第二相位令 $\gamma:=\pi-\alpha_2\in(0,\pi/2)$, 由
$O(\pi-\gamma)=\arctan(q\tan\gamma)$,
\begin{equation}
	c=c_2(\gamma,q):=\frac{\arctan(q\tan\gamma)}{\pi-\gamma},\qquad
	J_2^{(2)}(\gamma,q):=J(\pi-\gamma;c_2(\gamma,q)).
	\label{eq:c2}
\end{equation}

\begin{lemma}[真实相位曲线包含于二维盒]\label{lem:inclusion}
对 $(q,c)\in Q$, 有
\begin{equation}
	0.841<\alpha_1(q,c)<1.1220,\qquad
	0.655<\gamma(q,c)<1.0472.
\end{equation}
\end{lemma}

\begin{proof}
写 $F_1(q,c,x)=cx-\arctan(1/(q\tan x))$ 与
$F_2(q,c,\gamma)=c(\pi-\gamma)-\arctan(q\tan\gamma)$。对 $F_1$,
\begin{equation}
	\partial_xF_1=c+\frac q{\Phi_q(x)}>0,\qquad
	\partial_qF_1=\frac{\tan x}{1+q^2\tan^2x}>0,\qquad
	\partial_cF_1=x>0,
\end{equation}
故 $\alpha_1$ 关于 $q$ 与 $c$ 均严格递减; 对 $F_2$,
\begin{equation}
	\partial_\gamma F_2=-c-\frac{q\sec^2\gamma}{1+q^2\tan^2\gamma}<0,\qquad
	\partial_qF_2=-\frac{\tan\gamma}{1+q^2\tan^2\gamma}<0,\qquad
	\partial_cF_2=\pi-\gamma>0,
\end{equation}
故 $\gamma$ 关于 $q$ 严格递减、关于 $c$ 严格递增。于是
\begin{equation}
	\alpha_1(2,1/2)\le\alpha_1(q,c)\le\alpha_1(1,2/5),\qquad
	\gamma(2,2/5)\le\gamma(q,c)\le\gamma(1,1/2).
\end{equation}
端点闭式: 由 $\arctan(1/\tan x)=\pi/2-x$ 得 $\alpha_1(1,2/5)=5\pi/14$ 与
$\alpha_1(2,1/2)=\arccos(2/3)$ (后者因 $x/2=\arctan(1/(2\tan x))$ 等价于
$2\tan x\tan(x/2)=1$, 即 $\cos x=2/3$); 由 $\arctan(\tan\gamma)=\gamma$ 得
$\gamma(1,1/2)=\pi/3$。剩下的只是有理端点界:
\begin{itemize}
	\item $\arccos(2/3)>0.841$: 对 $x=841/1000$, 交错级数下界
	$\cos x\ge1-x^2/2+x^4/24-x^6/720>2/3$;
	\item $5\pi/14<1.1220$ 与 $\pi/3<1.0472$: 均只需 $\pi<3.1416$;
	\item $\gamma_*:=\gamma(2,2/5)>0.655$: 函数
	$h(\gamma)=\frac25(\pi-\gamma)-\arctan(2\tan\gamma)$ 严格递减且
	$h(\gamma_*)=0$, 而 $h(0.655)>0$: 由 $\pi>3.14159$ 得
	$\frac25(\pi-0.655)>0.9946$; 由 $\sin x\le x-x^3/6+x^5/120$ 与
	$\cos x\ge1-x^2/2+x^4/24-x^6/720$ 在 $x=131/200$ 处得
	$\tan(0.655)<0.7682$, 故 $2\tan(0.655)<1.5364$; 又
	$1/1.5364>0.6508$, 且交错级数下界
	$\arctan z\ge z-z^3/3+z^5/5-z^7/7+z^9/9-z^{11}/11$ 在
	$z=6508/10000$ 处给出 $\arctan(0.6508)>0.5767$, 于是
	\begin{equation}
		\arctan(2\tan(0.655))<\arctan(1.5364)<\frac\pi2-0.5767<0.9941<0.9946.
	\end{equation}
	故 $h(0.655)>0$, 即 $\gamma_*>0.655$。
\end{itemize}
\end{proof}

\subsubsection{\texorpdfstring{伴随命题 $(LOG)$ 的解析证明 (E1)}{伴随命题 (LOG) 的解析证明 (E1)}}\label{ss:log}

记 $G_k:=G(\alpha_k(c);c)$。伴随命题
\begin{equation}
	\frac{d}{dc}\log\frac{\widetilde M_{f1}}{\widetilde M_{f2}}=G_1-G_2<0
	\label{eq:LOG}
\end{equation}
由定理 \ref{thm:LOG} 完全解析证明 (E1), 取代旧 128 叶盒证书 (第 \ref{sec:certs}
节)。

\medskip
\noindent\textbf{第二相位的恒等式.} 置 $\gamma:=\pi-\alpha_2\in(0,\pi/2)$,
$A:=\pi-\gamma$ 与 $c=c_2(\gamma,q)$,
\begin{equation}
	\Phi:=\cos^2\gamma+q^2\sin^2\gamma,\qquad D:=q+c\Phi,\qquad
	W_0:=3-2A\cot\gamma,
\end{equation}
\begin{equation}
	P:=\frac{cA\Phi(q^2-1)\sin\gamma\cos\gamma}{D^2}.
\end{equation}
代入 $x=\pi-\gamma$ 于 \eqref{eq:G}, 由 $x\cot x=-A\cot\gamma$ 与
$\sin x\cos x=-\sin\gamma\cos\gamma$ 得精确恒等式
\begin{equation}
	G_2=-\frac{\Phi W_0}{D}-2P.
	\label{eq:G2id}
\end{equation}

\begin{lemma}\label{lem:G2m2}
设 $0.655\le\gamma\le\pi/3$, $1\le q\le2$, $c=c_2(\gamma,q)\in[0.4,0.5]$。
则 $G_2>-2$。
\end{lemma}

\begin{proof}
由 \eqref{eq:G2id}, 当 $W_0\ge0$ 时 $-\Phi W_0/D\ge-(\Phi/D)W_0$, 当
$W_0<0$ 时 $-\Phi W_0/D>0$; 因此
\begin{equation}
	G_2\ge-\frac{\Phi}{D}\max(W_0,0)-2P.
	\label{eq:G2low}
\end{equation}
下面给出三个估计。

\emph{(i) $\Phi/D<1$.} 记 $u:=\Phi/q=q\sin^2\gamma+\cos^2\gamma/q$。因
$\partial_q^2u=2\cos^2\gamma/q^3>0$, $u$ 在 $q\in[1,2]$ 上凸, 最大值在端点:
$u\le\max\{1,\,2\sin^2\gamma+\cos^2\gamma/2\}$。由 $\gamma\le\pi/3$,
\begin{equation}
	2\sin^2\gamma+\frac{\cos^2\gamma}{2}=\frac32\sin^2\gamma+\frac12
	\le\frac32\cdot\frac34+\frac12=\frac{13}{8},
\end{equation}
故 $u\le13/8$。由 $c\ge0.4$ 与 $t\mapsto t/(1+2t/5)$ 在 $t>0$ 上递增,
\begin{equation}
	\frac{\Phi}{D}=\frac{u}{1+cu}\le\frac{u}{1+\frac25u}
	\le\frac{13/8}{1+13/20}=\frac{65}{66}<1.
	\label{eq:PhiD}
\end{equation}
\emph{(ii) $W_0\le3-4\pi/(3\sqrt3)<0.582$.} 因
$\frac{d}{d\gamma}\bigl[(\pi-\gamma)\cot\gamma\bigr]
=-\cot\gamma-(\pi-\gamma)\csc^2\gamma<0$, 函数
$W_0=3-2(\pi-\gamma)\cot\gamma$ 在 $[0.655,\pi/3]$ 上严格递增, 故
$W_0\le W_0(\pi/3)=3-4\pi/(3\sqrt3)$。由 $\pi>3.1415$ 与 $\sqrt3<1.7321$,
\begin{equation}
	\frac{4\pi}{3\sqrt3}>\frac{4\cdot3.1415}{3\cdot1.7321}
	>\frac{12.566}{5.1963}>2.418,
\end{equation}
故 $W_0<3-2.418=0.582$。
\emph{(iii) $P<0.576$.} 由 $1\le\Phi\le q^2$, 函数
$f(\Phi):=\Phi(q^2-1)/(q+c\Phi)^2$ 满足
\begin{equation}
	f'(\Phi)=\frac{(q^2-1)(q-c\Phi)}{(q+c\Phi)^3}\ge0
	\qquad(\text{因 } q-c\Phi\ge q(1-cq)\ge0),
\end{equation}
故 $f(\Phi)\le f(q^2)=(q^2-1)/(1+cq)^2\le(q^2-1)/(1+0.4q)^2$; 末式
$h(q):=(q^2-1)/(1+0.4q)^2$ 满足 $h'(q)=2(q+0.4)/(1+0.4q)^3>0$, 故
$h(q)\le h(2)=25/27$。于是
\begin{equation}
	P=A\cdot(c\sin\gamma\cos\gamma)\cdot\frac{\Phi(q^2-1)}{D^2}
	\le A\cdot\frac14\cdot\frac{25}{27}\le\frac{25(\pi-0.655)}{108}<0.576,
\end{equation}
其中用 $c\sin\gamma\cos\gamma\le c/2\le1/4$, $A=\pi-\gamma\le\pi-0.655$,
且 $25(\pi-0.655)/108<0.576$ 由 $\pi<3.1416$ 得出。

组合 (i)--(iii) 于 \eqref{eq:G2low}: 由 $W_0<0.582$ 得
$\max(W_0,0)<0.582$, 结合 \eqref{eq:PhiD} 与 $P<0.576$,
\begin{equation}
	G_2>-(0.582)-2(0.576)=-1.734>-2.
	\label{eq:G2m2}
\end{equation}
\end{proof}

\begin{theorem}[$(LOG)$]\label{thm:LOG}
对一切 $q>1$ 与 $0<c<1/2$,
\begin{equation}
	G_1-G_2<0.
\end{equation}
\end{theorem}

\begin{proof}
若 $G_2\ge0$, 引理 \ref{lem:G1} 给出 $G_1<0$, 故 $G_1-G_2<0$。若 $G_2<0$,
则由 \eqref{eq:B} 得 $(q,c)\in(1,2)\times(0.4,0.5)$; 引理 \ref{lem:inclusion}
及其证明给出 $0.655<\gamma\le\pi/3$ 与 $c=c_2(\gamma,q)\in(0.4,0.5)$, 故引理
\ref{lem:G2m2} 给出 $G_2>-2$。另一方面, 定理 \ref{thm:j1e1} 证明的第 (iii)
步在 $T_1$ 上证明 $G<-2$, 而由 \eqref{eq:T1} 与引理 \ref{lem:inclusion}
该步适用于 $\alpha_1$, 故 $G_1<-2$。于是 $G_1-G_2<-2-G_2<0$。两种情形穷尽
全部参数。
\end{proof}

定理 \ref{thm:j1e1} (纯 E1) 给出 $J_1^{(2)}\ge6499/7500>0$ 于 $T_1$ 的闭包;
定理 \ref{thm:j2e1} (纯 E1, 见第 \ref{ss:j2e1} 子节) 给出
\begin{equation}
	J_2^{(2)}(\gamma,q)<0\quad\text{于}\ [0.655,1.0472]\times[1,2].
	\label{eq:Jcert}
\end{equation}
结合引理 \ref{lem:inclusion} 与 $\widetilde M_{f1},\widetilde M_{f2}>0$,
由 \eqref{eq:Fep2raw} 得
\begin{equation}
	\widetilde F_e''(q,c)>0\quad\text{于}\ Q.
	\label{eq:Fep2}
\end{equation}
在端点 $c=1/2$, 令
\begin{equation}
	x=2\arcsin\frac{1}{\sqrt{2(q+1)}}\in(0,\pi/3).
\end{equation}
精确公式为
\begin{equation}
	\widetilde F_e'(q,1/2)=\frac{2\pi(\cos x-1)^3}{\sin^3x}P(x),
\end{equation}
\begin{equation}
	P(x)=3x^2+6x\sin x-3\pi x-3\pi\sin x+\pi^2.
\end{equation}
又
\begin{equation}
	P(x)=(\pi-3x)^2+3(x-\sin x)(\pi-2x)>0.
\end{equation}
由于 $(\cos x-1)^3<0$ 且 $\sin^3x>0$,
\begin{equation}
	\widetilde F_e'(q,1/2)<0.
	\label{eq:Fep12}
\end{equation}

若 $G_2<0$, 则由 \eqref{eq:B} 点 $(q,c)$ 位于 $Q$。对固定 $q$, \eqref{eq:Fep2}
说明 $\widetilde F_e'$ 随 $c$ 严格增加; 因 $c<1/2$, 结合 \eqref{eq:Fep12} 得
\begin{equation}
	\widetilde F_e'(q,c)<\widetilde F_e'(q,1/2)<0.
\end{equation}
若 $G_2\ge0$, 则已由引理 \ref{lem:G1} 和 \eqref{eq:Fep} 得到同一结论。两种情形
穷尽全部参数, 定理 \ref{thm:keylemma} 得证。


\subsubsection{证据分层与真曲线区域上的解析化路线}

按第 \ref{rem:evidence} 节的约定, 本子节的证据分层如下: 引理 \ref{lem:inclusion}
的端点界、端点值 \eqref{eq:Fep12}, 以及下面的恒等式与全部单变量估计都是严格
解析证明 (E1); 定理 \ref{thm:j1e1} 以完全初等的方式证明
$J_1^{(2)}\ge6499/7500>0$ 于 $T_1$ 上; \eqref{eq:Jcert} 中 $J_2^{(2)}$ 的
不等式由定理 \ref{thm:j2e1} 完全解析证明 (E1), 其单变量事实由注 \ref{rem:env}
的有理包络方法给出 (证书总表见附录 \ref{app:cert}); 注 \ref{rem:explore} 的范围数据
是 E3 扫描, 仅作侦察与
交叉检验, 不作为任何结论依据。

对第一相位, 记
\begin{equation}
	T_1:=\{(x,q): 0.841<x<1.1220,\ 1<q<2,\ 0.4<c_1(x,q)<0.5\}.
	\label{eq:T1}
\end{equation}
由 $c_1(\alpha_1(q,c),q)=c$ 与引理 \ref{lem:inclusion}, $Q$ 上真实相位曲线的像
包含于 $T_1$; 第二相位同理,
\begin{equation}
	T_2:=\{(\gamma,q): 0.655<\gamma<1.0472,\ 1<q<2,\ 0.4<c_2(\gamma,q)<0.5\}.
	\label{eq:T2}
\end{equation}
因此只需在 $T_1,T_2$ 上证明 \eqref{eq:Jcert} 型不等式; 对 $J_1^{(2)}$ 由定理
\ref{thm:j1e1}, 对 $J_2^{(2)}$ 由定理 \ref{thm:j2e1}, 均已完全解析化, 无需
证书盒。

以下恒等式均为 (E1)。记
\begin{equation}
	u:=\frac{x\Phi_q(x)}{q+c\Phi_q(x)},\qquad
	A:=\frac3x+2\cot x,\qquad
	H:=\frac{2c(q^2-1)\sin x\cos x}{q+c\Phi_q(x)},
\end{equation}
则 $G=u(H-A)$, 且沿真实曲线 $c=c_1(x,q)$ 有分解
\begin{equation}
	J_1^{(2)}=G^2+G_c-\frac{x\Phi_q(x)}{q+c\Phi_q(x)}G_x,
	\label{eq:jdec}
\end{equation}
其中 $G_c:=\partial_cG$, $G_x:=\partial_xG$ 是 $G$ 在固定 $(q,c)$ 下的偏导数
(与 \eqref{eq:J} 中记号一致, 不是沿相位曲线的导数); $J_2^{(2)}$ 有同构分解。

\begin{theorem}[$J_1^{(2)}$ 的完全解析下界]\label{thm:j1e1}
在 $T_1$ 的闭包上,
\begin{equation}
	J_1^{(2)}(x,q)\ge\frac{6499}{7500}>\frac{1733}{2000}>0.
	\label{eq:j1bound}
\end{equation}
\end{theorem}

\begin{proof}
记 $D:=q+c\Phi_q(x)$, $\Phi:=\Phi_q(x)$, $W:=3+2x\cot x$。证明由七个初等
步骤组成。
\emph{(i) $\Phi/D\ge2/3$.} 因 $x>\pi/4$, 有 $\cos2x<0$; 而
\begin{equation}
	\Phi-q=q^2\sin^2x-q+\cos^2x,\qquad
	\frac{d}{dq}(\Phi-q)=2q\sin^2x-1\ge2\sin^2x-1=-\cos2x>0,
\end{equation}
故 $\Phi\ge q$ 对 $q\ge1$ 成立, 于是 $q/\Phi\le1$; 结合 $c\le1/2$,
\begin{equation}
	\frac{\Phi}{D}=\frac1{q/\Phi+c}\ge\frac1{1+1/2}=\frac23.
	\label{eq:phid23}
\end{equation}
\emph{(ii) $u\ge2x/3$ 且 $u_x\ge2/3$.} 由 \eqref{eq:phid23} 得
$u=x\Phi/D\ge2x/3$。又
\begin{equation}
	u_x=\frac{\Phi D+2xq(q^2-1)\sin x\cos x}{D^2}
	=\frac{\Phi}{D}+\frac{2xq(q^2-1)\sin x\cos x}{D^2}\ge\frac{\Phi}{D}\ge\frac23.
	\label{eq:ux}
\end{equation}
\emph{(iii) $G<-2$.} 由
\begin{equation}
	D-c(q^2-1)\sin^2x=q+c\bigl(\cos^2x+q^2\sin^2x-(q^2-1)\sin^2x\bigr)=q+c>0
\end{equation}
得 $H<2\cot x$, 故
\begin{equation}
	G=u(H-A)<u\Bigl(2\cot x-\frac3x-2\cot x\Bigr)=-\frac{3\Phi}{D}\le-2,
\end{equation}
从而 $G^2\ge4$。

\emph{(iv) $G_c\ge187/100$.} 分解 $G_c=t_1+t_2$, 其中
\begin{equation}
	t_1:=\frac{\Phi^2W}{D^2},\qquad
	t_2:=\frac{2x\Phi(q^2-1)\sin x\cos x\,(q-c\Phi)}{D^3}.
\end{equation}
由 $\Phi\le q^2$ 与 $c\le1/2$, $q\le2$ 得 $q-c\Phi\ge q-cq^2=q(1-cq)\ge0$,
故 $t_2\ge0$。对 $t_1$ 分两段。若 $x\le\pi/3$, 因 $x\cot x$ 在 $(0,\pi/2)$
递减 ($\frac{d}{dx}x\cot x=(\sin x\cos x-x)/\sin^2x<0$), 结合 \eqref{eq:phid23},
\begin{equation}
	t_1\ge\Bigl(\frac23\Bigr)^2\Bigl(3+\frac{2\pi}{3\sqrt3}\Bigr)
	=\frac43+\frac{8\pi}{27\sqrt3}>\frac43+\frac{161}{300}=\frac{187}{100},
\end{equation}
其中末式用 $\pi>3.1415$, $\sqrt3<1.7321$ 使 $8\pi/(27\sqrt3)>
8\cdot3.1415/(27\cdot1.7321)>161/300$。若 $x\ge\pi/3$, 则沿曲线
$c=c_1(x,q)$, 由 $dc_1/dq<0$ 与 $\partial_c(\Phi/D)=-\Phi^2/D^2$ 得
\begin{equation}
	\frac{d}{dq}\frac{\Phi}{D}
	=\frac{q^2\sin^2x-\cos^2x}{D^2}
	+\frac{\Phi^2\tan x}{xD^2(q^2\tan^2x+1)}>0
\end{equation}
(第一项非负因 $x>\pi/4$), 于是 $t_1(x,q)\ge t_1(x,1)=(2x/\pi)^2W(x)=:f(x)$。
而
\begin{equation}
	f'(x)=\frac{8x}{\pi^2}\bigl(3+3x\cot x-x^2\csc^2x\bigr),
\end{equation}
且对 $x\in[\pi/3,5\pi/14]$, 由 $x\cot x$ 递减及 $x\le5\pi/14$, $\sin x\ge\sin(\pi/3)$,
\begin{equation}
	x\cot x\ge\frac{5\pi}{14}\tan\frac{\pi}{7},\qquad
	x^2\csc^2x\le\frac{4}{3}\Bigl(\frac{5\pi}{14}\Bigr)^2,
\end{equation}
故 $3+3x\cot x-x^2\csc^2x\ge3+3(5\pi/14)\tan(\pi/7)-\frac43(5\pi/14)^2>0$
(用 $\tan t\ge t+t^3/3+2t^5/15$ 于 $t=\pi/7$ 与 $\pi\in(3.1415,3.1416)$ 的有理界)。

对 $x\in[5\pi/14,1122/1000]$ (区间非空因 $\pi<3.1416$, 且含于 $(0,\pi/2)$ 因
$\pi>3.14$), 由 $x\cot x$ 递减、$x/\sin x$ 递增
($\frac{d}{dx}(x/\sin x)=(\sin x-x\cos x)/\sin^2x>0$), 以及引理
\ref{lem:envseries} 在有理点 $x_0:=1122/1000$ 的交错级数夹逼
\begin{equation}
	\frac{9009}{10000}<\sin x_0<\frac{9010}{10000},\qquad
	\frac{4338}{10000}<\cos x_0<\frac{4340}{10000},
	\label{eq:env1122}
\end{equation}
(四条均为有限次精确有理运算, 对应部分和 $S_3,S_4$ 与 $C_3,C_4$), 得
\begin{equation}
	x\cot x\ge x_0\frac{\cos x_0}{\sin x_0}
	\ge\frac{1122}{1000}\cdot\frac{4338}{9010}
	>\frac{1122}{1000}\cdot\frac{48}{100},\qquad
	\frac{x}{\sin x}\le\frac{x_0}{\sin x_0}
	<\frac{1122/1000}{9009/10000}=\frac{11220}{9009}<\frac{1246}{1000},
\end{equation}
其中 $4338/9010>48/100$ 与 $11220/9009<1246/1000$ 是直接有理不等式。于是
\begin{equation}
	3+3x\cot x-x^2\csc^2x\ge
	3+3\cdot\frac{1122}{1000}\cdot\frac{48}{100}-\Bigl(\frac{1246}{1000}\Bigr)^2
	=\frac{765791}{250000}>0.
\end{equation}
故 $f$ 在 $[5\pi/14,1122/1000]$ 递增; 与前半段在 $5\pi/14$ 处拼合, $f$ 在
$[\pi/3,1122/1000]$ 递增, $t_1\ge f(x)\ge f(\pi/3)=4/3+8\pi/(27\sqrt3)
>187/100$。

\emph{(v) $C:=3/x^2+2\csc^2x\le8$.} $C$ 在 $(0,\pi/2)$ 递减, 且
$\sin(841/1000)\ge841/1000-(841/1000)^3/6>0.7418$, 故
$\sin^2(841/1000)>0.55$; 于是
\begin{equation}
	C\le\frac3{(841/1000)^2}+\frac2{\sin^2(841/1000)}
	<\frac{3000}{707}+\frac{200}{55}<8.
\end{equation}

\emph{(vi) $u\le89/100$.} 记 $\theta:=c_1(x,q)x$ 与
$W_0:=x\sin\theta\cos\theta+\theta\sin x\cos x$。恒等式
$\Phi/D=x\sin x\cos x/W_0$ 给出 $u=x^2\sin x\cos x/W_0$。定义
\begin{equation}
	u_{\rm c}(x):=\frac{x\sin2x}{\sin(4x/5)+0.4\sin2x}.
\end{equation}
因 $c_1>0.4$, 有 $2\theta\ge4x/5$; 又 $2\theta\le x<1.122<\pi/2$, 故
$\sin2\theta\ge\sin(4x/5)$。由
$u\le u_{\rm c}\iff x\sin(4x/5)+0.8x\sin x\cos x\le2W_0$ 及
$2W_0\ge x\sin2\theta+0.8x\sin x\cos x$, 得 $u\le u_{\rm c}(x)$。其次,
$u_{\rm c}(x)\le89/100$ 等价于 $F(x)\ge0$, 其中
\begin{equation}
	F(x):=\frac{89}{100}\sin\frac{4x}{5}-\Bigl(x-\frac{89}{250}\Bigr)\sin2x.
\end{equation}
计算得
\begin{equation}
	F''(x)=-\frac{356}{625}\sin\frac{4x}{5}-4\cos2x+4\Bigl(x-\frac{89}{250}\Bigr)\sin2x.
\end{equation}
令 $y:=2x\in[841/500,1122/500]\subset(\pi/2,\pi)$ 与
$g(y):=(y/2-89/250)\sin y-\cos y$, 则 $F''=-(356/625)\sin(4x/5)+4g(y)$,
\begin{equation}
	g'(y)=\frac32\sin y+\Bigl(\frac y2-\frac{89}{250}\Bigr)\cos y
	\ge\frac32\sin(0.8975)-0.766\cos(0.8975)>0,
\end{equation}
故 $g$ 递增, $g(y)\ge g(841/500)=0.485\sin(841/500)-\cos(841/500)$; 用
$\sin(841/500)\ge\sin(17/10)\ge\cos(13/100)$ 与
$-\cos(841/500)=\cos(\pi-841/500)\ge\cos(14596/10000)$, 以及交错级数在有理点
的包络 ($\sin(8976/10000)\le0.78193$, $\cos(13/100)\ge0.99155$,
$\cos(14596/10000)\ge0.11047$, 全部由复现脚本逐项核对, 见第 \ref{sec:certs}
节内容哈希 L6), 得
\begin{equation}
	F''(x)\ge4\Bigl(0.485\cos\frac{13}{100}+\cos\frac{14596}{10000}\Bigr)
	-\frac{356}{625}\cdot0.78193>\frac32>0.
\end{equation}
同一组有理包络给出
\begin{equation}
	F'\Bigl(\frac{24}{25}\Bigr)\in\Bigl(-\frac1{20},0\Bigr),\qquad
	F'\Bigl(\frac{97}{100}\Bigr)>0,\qquad
	F\Bigl(\frac{24}{25}\Bigr)\ge\frac{49}{1000},\qquad
	F\Bigl(\frac{97}{100}\Bigr)\ge\frac{49}{1000}.
\end{equation}
由 $F''>0$ 知 $F'$ 严格递增, 故 $F$ 在 $[841/1000,1122/1000]$ 上先减后增,
唯一极小点位于 $(24/25,97/100)$ 内, 于是 $F\ge49/1000>0$, 即
$u_{\rm c}\le89/100$, 从而 $u\le89/100$。

\emph{(vii) 组合.} 由 $H_x<0$ (因 $\cos2x<0$)、$H-A<-3/x$ 与
$G_x=u_x(H-A)+u(H_x-A_x)$, $A_x=-3/x^2-2\csc^2x=-C$, 得
\begin{equation}
	uG_x=u\,u_x(H-A)+u^2H_x+u^2C
	\le u\,u_x\Bigl(-\frac3x\Bigr)+u^2C
	=-\frac{3uu_x}{x}+u^2C.
\end{equation}
由 (ii) 得 $3uu_x/x\ge4/3$, 由 (v)(vi) 得 $u^2C\le(89/100)^2\cdot8$。于是
\begin{equation}
	J_1^{(2)}=G^2+G_c-uG_x
	\ge4+\frac{187}{100}-\Bigl(\Bigl(\frac{89}{100}\Bigr)^2\cdot8-\frac43\Bigr)
	=\frac{6499}{7500},
\end{equation}
即 \eqref{eq:j1bound}。
\end{proof}

单变量 $q=1$ 情形可以完全解析 (E1)。由 $\Phi_q\equiv1$、
$c_1(x,1)=(\pi/2-x)/x$ 与 $D=\pi/(2x)$,
\begin{equation}
	J_1^{(2)}(x,1)=\Bigl(\frac{2x}{\pi}\Bigr)^2N(x),\qquad
	N(x)=12+16x\cot x+2x^2\cot^2x-2x^2.
	\label{eq:j1q1}
\end{equation}
对 $x\in[\pi/3,5\pi/14]$, 记 $z:=x\cot x$。由 $\cot x\ge\cot(5\pi/14)=\tan(\pi/7)$
与 $\tan t\ge t+t^3/3$ ($t\ge0$) 得 $z\ge(\pi/3)\tan(\pi/7)>1/2$; 又
$z\le x<1.13$。$z\mapsto z^2+8z+6$ 在 $(0,\infty)$ 递增, 且
$x^2\le(5\pi/14)^2<1.26$, 于是
\begin{equation}
	N(x)=2(z^2+8z+6)-2x^2
	\ge2\Bigl(\frac14+4+6\Bigr)-2\cdot1.26>17.9>0.
\end{equation}
故 $J_1^{(2)}(x,1)>0$ 于 $x\in[\pi/3,5\pi/14]$; 角点
$(x,q)=(\pi/3,1)$ 处 $J_1^{(2)}=4N(\pi/3)/9\approx8.98$。

对第二相位, 令 $x:=\pi-\gamma\in[2\pi/3,5\pi/7]$, 则 $c_2(\gamma,1)=\gamma/(\pi-\gamma)$
与 $D=\pi/x$, 且 $J_2^{(2)}(\gamma,1)=x^2N(x)/\pi^2$。记
$z:=-x\cot x=(\pi-\gamma)\cot\gamma$。由 $\gamma\in[2\pi/7,\pi/3]$ 得
\begin{equation}
	z\in\Bigl[\frac{2\pi}{3\sqrt3},\,\frac{5\pi}{7}\cot\frac{2\pi}{7}\Bigr]
	\subset(1.19,2.5),
\end{equation}
其中 $z\mapsto z^2-8z+6$ 递减 (因 $z<4$)。用 $\sqrt3<7/4$ 得
$z\ge2\pi/(3\sqrt3)>8\pi/21$, 用 $x\ge2\pi/3$, 于是
\begin{equation}
	N(x)=2(z^2-8z+6)-2x^2
	\le2\Bigl(4-8\cdot\frac{8\pi}{21}+6\Bigr)-2\Bigl(\frac{2\pi}{3}\Bigr)^2<-7<0.
\end{equation}
故 $J_2^{(2)}(\gamma,1)<0$ 于 $\gamma\in[2\pi/7,\pi/3]$; 角点
$(\gamma,q)=(\pi/3,1)$ 处 $J_2^{(2)}=4N(2\pi/3)/9\approx-5.87$。

对第一相位, 定理 \ref{thm:j1e1} 已把 $J_1^{(2)}$ 的正性完全解析化。
对第二相位, 完全解析化由下面的定理 \ref{thm:j2e1} 完成, 取代原二维叶盒证书
与引理 (M1$'$)--(M3$'$) 路线。

\subsubsection{\texorpdfstring{$J_2^{(2)}$ 的完全解析证明 (E1)}{J2 的完全解析证明 (E1)}}\label{ss:j2e1}

记
\begin{equation}
	\gamma\in[0.655,1.0472],\quad q\in[1,2],\quad
	t:=\arctan(q\tan\gamma),\quad A:=\pi-\gamma,
\end{equation}
\begin{equation}
	\mathrm{sg}:=\sin\gamma,\quad \mathrm{cg}:=\cos\gamma,\quad
	\mathrm{st}:=\sin t,\quad \mathrm{ct}:=\cos t,\quad
	D:=\sqrt{1+3\,\mathrm{sg}^2},\quad
	\tau:=\arctan(2\tan\gamma).
\end{equation}
(记号 $\mathrm{sg},\mathrm{cg},\mathrm{st},\mathrm{ct}$ 仅在本子节使用。)

\begin{lemma}[代数分解]\label{lem:j2dec}
沿真实曲线 $c_2(\gamma,q)=t/A$,
\begin{equation}
	J_2^{(2)}(\gamma,q)=\frac{N}{16\Delta^4},\qquad
	\Delta:=A\,\mathrm{st}\,\mathrm{ct}+t\,\mathrm{sg}\,\mathrm{cg}>0,
	\label{eq:j2delta}
\end{equation}
其中 $N$ 是 $A,t,\mathrm{sg},\mathrm{cg},\mathrm{st},\mathrm{ct}$ 的整系数
多项式, 且
\begin{equation}
	N=32A^2\,\mathrm{cg}\,W,\qquad W:=\mathcal W_1+\cdots+\mathcal W_8,
	\label{eq:wdec}
\end{equation}
\begin{align}
	\mathcal W_1&=-2A^3B_1\,\mathrm{st}^2\,\mathrm{ct}^4,
	&\mathcal W_2&=A^2\,\mathrm{cg}\,B_2\,\mathrm{st}^2\,\mathrm{ct}^2,\nonumber\\
	\mathcal W_3&=-2A^3\,\mathrm{sg}\,t\,\mathrm{st}\,\mathrm{ct}^5,
	&\mathcal W_4&=A^2\,\mathrm{sg}\,t\,B_4\,\mathrm{st}\,\mathrm{ct}^3,\nonumber\\
	\mathcal W_5&=-A\,\mathrm{cg}^2\,\mathrm{sg}\,t\,B_5\,\mathrm{st}\,\mathrm{ct},
	&\mathcal W_6&=4A^2\,\mathrm{cg}\,\mathrm{sg}^2\,t^2\,\mathrm{ct}^4,\nonumber\\
	\mathcal W_7&=-A\,\mathrm{cg}\,\mathrm{sg}^2\,t^2\,B_7\,\mathrm{ct}^2,
	&\mathcal W_8&=6\,\mathrm{cg}^3\,\mathrm{sg}^4\,t^2,
	\label{eq:wterms}
\end{align}
其中括号因子为
\begin{equation}
	\begin{aligned}
		B_1&:=A\,\mathrm{cg}-2\,\mathrm{sg},&
		B_2&:=4A^2\,\mathrm{cg}^2-A^2-12A\,\mathrm{cg}\,\mathrm{sg}+6\,\mathrm{sg}^2,\\
		M&:=2A^2\,\mathrm{cg}^2-A^2-8A\,\mathrm{cg}\,\mathrm{sg}+6\,\mathrm{sg}^2
		=B_2-2A\,\mathrm{cg}\,B_1,\\
		B_4&:=7A\,\mathrm{cg}^2-A\,\mathrm{sg}^2-4\,\mathrm{cg}\,\mathrm{sg},\\
		B_5&:=A^2\,\mathrm{cg}^2-A^2\,\mathrm{sg}^2+2A^2+12A\,\mathrm{cg}\,\mathrm{sg}-12\,\mathrm{sg}^2,\\
		B_7&:=3A\,\mathrm{cg}^2+A\,\mathrm{sg}^2+8\,\mathrm{cg}\,\mathrm{sg},\qquad
		G_5:=B_5-AB_4.
	\end{aligned}
	\label{eq:brackets}
\end{equation}
\end{lemma}

\begin{proof}
\eqref{eq:j2delta} 由 $J(\pi-\gamma;c_2)$ 的定义直接展开 (把相位方程显式代入后
通分); $N$ 的整系数多项式由符号计算生成 (misc/t3\_NJ2.json), 其显式形式不重要,
只需精确分解 \eqref{eq:wdec}。\eqref{eq:wdec}--\eqref{eq:wterms} 在模
$\mathrm{sg}^2+\mathrm{cg}^2=1$ 与 $\mathrm{st}^2+\mathrm{ct}^2=1$ 下是精确
多项式恒等式, 由符号计算 (misc/\_verify\_identity.py 与
misc/zz\_rebuild\_check1.py) 精确验证, 不是数值近似。$\Delta>0$ 因各项为正;
$M=B_2-2A\,\mathrm{cg}\,B_1$ 直接展开即得。
\end{proof}

\begin{lemma}[轨迹几何]\label{lem:track}
在 $[0.655,1.0472]\times[1,2]$ 上:
\begin{enumerate}
	\item[(i)] $t\in[\gamma,\tau]$, 且对 $z:=\mathrm{ct}^2$ 有 $z\in[z_-,z_+]$,
	其中 $z_-:=\mathrm{cg}^2/D^2$, $z_+:=\mathrm{cg}^2$;
	\item[(ii)] $z_-+z_+\le1$, 因此 $z(1-z)\ge z_-(1-z_-)=4\,\mathrm{sg}^2\,\mathrm{cg}^2/D^4$
	对一切 $z\in[z_-,z_+]$ 成立;
	\item[(iii)] $t\,\mathrm{st}\,\mathrm{ct}^5\ge2\tau\,\mathrm{sg}\,\mathrm{cg}^5/D^6$;
	\item[(iv)] $\frac t2\sin(2t)\ge m:=\frac{3164}{10000}$;
	\item[(v)] $D\le2$.
\end{enumerate}
\end{lemma}

\begin{proof}
(i) $\arctan$ 严格递增且 $1\le q\le2$, 故 $t\in[\gamma,\tau]\subset(0,\pi/2)$。
$\cos$ 在 $(0,\pi/2)$ 递减, 故 $\mathrm{ct}\in[\cos\tau,\cos\gamma]$, 而
\begin{equation}
	\cos\tau=\frac1{\sqrt{1+4\tan^2\gamma}}
	=\frac{\mathrm{cg}}{\sqrt{1+3\,\mathrm{sg}^2}}=\frac{\mathrm{cg}}D.
\end{equation}
(ii) $z_-+z_+\le1\iff \mathrm{cg}^2(1+1/D^2)\le1\iff 3u^2+2u-1\ge0$
($u:=\mathrm{sg}^2$) $\iff u\ge1/3$; 而 $\sin(0.655)\ge0.655-0.655^3/6>0.6>1/\sqrt3$,
故 $u\ge\sin^2(0.655)>1/3$。于是对 $z\in[z_-,z_+]$,
\begin{equation}
	z(1-z)-z_-(1-z_-)=(z-z_-)(1-z-z_-)\ge0,
\end{equation}
且 $z_-(1-z_-)=\mathrm{cg}^2(D^2-\mathrm{cg}^2)/D^4=4\,\mathrm{sg}^2\,\mathrm{cg}^2/D^4$
(用 $D^2-\mathrm{cg}^2=1+3\,\mathrm{sg}^2-\mathrm{cg}^2=4\,\mathrm{sg}^2$)。
(iii) 令 $f(t):=t\sin t\cos^5t$。于 $t\in[0.655,\pi/2)$:
\begin{equation}
	f'(t)=\cos^4t\bigl(\sin t\cos t+t(1-6\sin^2t)\bigr)
	\le\cos^4t\Bigl(\frac12+0.655\cdot(1-6\sin^2(0.655))\Bigr)<0,
\end{equation}
因 $\sin t\cos t\le1/2$, $\sin t\ge\sin(0.655)>0.6$ 使 $1-6\sin^2t<-1.16$,
且 $t\ge0.655$。故 $f$ 递减, $f(t)\ge f(\tau)$; 由 $\sin\tau=2\,\mathrm{sg}/D$ 与
$\cos\tau=\mathrm{cg}/D$ 得 $f(\tau)=\tau(2\,\mathrm{sg}/D)(\mathrm{cg}/D)^5
=2\tau\,\mathrm{sg}\,\mathrm{cg}^5/D^6$。
(iv) 令 $h(t):=\frac t2\sin2t=t\sin t\cos t$。先证 $h$ 在 $[0.655,13/10]$ 凹:
$h''(t)=2(\cos2t-t\sin2t)$。于 $[0.655,\pi/4]$: $\cos2t\le\cos(2\cdot0.655)<26/100$
且 $t\sin2t\ge0.655\sin(2\cdot0.655)>3/5$ (交错级数在有理点 $131/100$ 的包络),
故 $h''\le2(26/100-3/5)<0$; 于 $[\pi/4,13/10]$: $13/10<\pi/2$, 故
$\cos2t\le0\le t\sin2t$, $h''\le0$ (除孤立点外严格)。又 $\tau(0.655)>\pi/4>0.655$ (因 $2\tan(0.655)>2\cdot0.655>1$) 且
$\tau(1.0472)<13/10$ (表 \ref{tab:envpoints}), 故
$[\gamma,\tau]\subset[0.655,13/10]$。凹函数在区间上的最小值在端点取得, 而
$h(0.655)\ge m$ 与 $h(13/10)\ge m$ (表 \ref{tab:envpoints}), 故 $h(t)\ge m$ 于
$[0.655,13/10]$, 特别地 $h(\gamma)\ge m$ 与 $h(\tau)\ge m$, 得 (iv)。
(v) $D^2=1+3\,\mathrm{sg}^2\le4$。
\end{proof}

定义
\begin{equation}
	\begin{gathered}
	T_{A,B_2}:=\frac{4A^2\,\mathrm{sg}^2\,\mathrm{cg}^4|B_2|}{D^4},\qquad
	T_{A,M}:=\frac{4A^2\,\mathrm{sg}^2\,\mathrm{cg}^4|M|}{D^4},\\
	T_A:=T_{A,B_2}\ \text{若}\ \gamma\le\gamma_0,\quad
	T_A:=T_{A,M}\ \text{若}\ \gamma\ge\gamma_0,
	\end{gathered}
\end{equation}
\begin{equation}
	T_B:=\frac{2A^3\,\mathrm{sg}^2\,\tau\,\mathrm{cg}^5}{D^5},\qquad
	T_C:=m\,G_5\,A\,\mathrm{sg}\,\mathrm{cg}^2,\qquad
	F:=\tau^2\,\mathrm{cg}\,\mathrm{sg}^2,
\end{equation}
\begin{equation}
	Q(z):=4A^2z^2-AB_7z+6\,\mathrm{cg}^2\,\mathrm{sg}^2,\qquad
	Q_-:=Q(z_-),\qquad Q_+:=Q(z_+),\qquad T_D:=F\max(Q_-,0).
\end{equation}

\begin{remark}[有理包络方法]\label{rem:env}
下面引理 \ref{lem:brackets} 与引理 \ref{lem:track}(iv) 中的全部单变量事实
(共 55 项) 由下列\emph{有理包络方法}逐项给出 E1 证明, 不再使用任何区间验证器
内核。方法的四个组成部分:
\begin{enumerate}
	\item[(a)] \emph{交错级数包络}: 对有理点 $x\in(0,3/2)$, $\sin x$ 与 $\cos x$
	的 Taylor 部分和交替夹逼真值, 余项被下一项控制 (引理 \ref{lem:envseries});
	对 $x\in(0,1]$, $\arctan x$ 同理, 对 $x>1$ 用
	$\arctan x=\pi/2-\arctan(1/x)$; $\pi$ 用 Machin 公式
	$\pi=16\arctan(1/5)-4\arctan(1/239)$ 的两个反正切项包络组合。本文对
	$\sin,\cos$ 取 12 项、对 $\arctan$ 取 22 项; 由此得到的原语包络认证宽度最坏
	$\approx1.8\times10^{-10}$ (在 $\tau(131/200)$ 行, 由 $\arctan$ 余项主导), 远
	小于证书最小裕量 ($\approx2.6\times10^{-5}$), 不影响任何结论; 细节见引理
	\ref{lem:envseries} 末尾。
	\item[(b)] \emph{单调包络}: 在 $[0.655,1.0472]$ 上 $\sin\gamma$ 增、
	$\cos\gamma$ 减、$A=\pi-\gamma$ 减、$D=\sqrt{1+3\sin^2\gamma}$ 增、
	$\tau=\arctan(2\tan\gamma)$ 增 (因 $\tau'=2/D^2>0$), 故任意小区间
	$[a,b]\subset[0.655,1.0472]$ 上这些原语的值域由端点处的 (a) 包络夹出。
	\item[(c)] \emph{精确有理区间算术}: 四则运算与正整数幂作用于有理端点区间时
	按区间自然定义 (端点组合取 min/max) 计算, 每步都是有限次精确有理数运算; 由
	(a)(b) 的原语包络经此算术得到任何代数组合 ($B_1,\dots,G_5$, $Q_\pm$, $F$,
	$T_{A,B_2}$, $T_{A,M}$, $T_B$, $T_C$ 等) 的认证区间。
	\item[(d)] \emph{泰勒模型}: 设 $f$ 在 $[a,b]$ 上 $C^2$, $c=(a+b)/2$,
	$w=(b-a)/2$, $M:=\sup_{[a,b]}|f''|$。由中值定理
	$f'(\gamma)=f'(c)+f''(\xi)(\gamma-c)$ 得
	$f'(\gamma)\in f'(c)+[-Mw,Mw]$; 若 $f'(c)$ 的包络 (由 (a)(c)) 与 $M$
	(由 (b)(c) 对 $f''$ 的链式法则表达式求包络) 使 $f'(c)$ 的下界 $-Mw>0$
	(或上界 $+Mw<0$), 则 $f'>0$ (或 $f'<0$) 于 $[a,b]$。值泰勒模型
	$f(\gamma)\in f(c)+[-M'w,M'w]$ ($M'=\sup|f'|$) 同理。判据的完整陈述见
	引理 \ref{lem:envtaylor}。
\end{enumerate}
表 \ref{tab:envprims}--\ref{tab:envderiv} (附录 \ref{app:cert}) 逐项列出由此
得到的证书: 每个区间都是一条有限精确有理数不等式链, 可人工复核; 生成脚本
misc/e1\_certgen.py 与认证台账 misc/e1\_cert\_ledger.json 只用于生成与复现,
其正确性不是结论的依据。
\end{remark}

\begin{lemma}[括号符号与单调性]\label{lem:brackets}
在 $[0.655,1.0472]$ 上:
\begin{enumerate}
	\item[(i)] $B_1$ 严格递减, $B_1(0.85)\ge1/200>0$ 且 $B_1(0.86)\le-1/50<0$;
	故存在唯一 $\gamma_0\in(0.85,0.86)$ 使 $B_1(\gamma_0)=0$;
	\item[(ii)] $B_2<0$, $M<0$, $B_4>0$, $G_5>0$;
	\item[(iii)] $Q_+<0$; $Q_-$ 严格递增, $Q_-(1.0014)\le-1/10000<0<Q_-(1.0472)\le33/200$;
	\item[(iv)] $F$ 在 $[1.0014,1.0472]$ 上严格递增, $F(1.0472)\le63/100$;
	\item[(v)] $T_{A,B_2}$ 递增于 $[0.655,0.72]$ 与 $[0.72,0.723]$,
	$\ge27/10$ 于 $[0.723,0.724]$, 递减于 $[0.724,0.73]$, $[0.73,0.85]$ 与 $[0.85,0.86]$;
	$T_{A,M}$ 递减于 $[0.85,0.86]$ 与 $[0.86,1.0472]$;
	\item[(vi)] $T_B$ 于 $[0.655,1.0472]$ 递减; $T_C$ 递增于 $[0.655,0.82]$,
	$\ge19/10$ 于 $[0.82,0.83]$, 递减于 $[0.83,1.0472]$。
\end{enumerate}
\end{lemma}

\begin{proof}
$B_1'(\gamma)=-3\cos\gamma-(\pi-\gamma)\sin\gamma<0$ (因
$\gamma\in(0,\pi/2)$), 得 (i) 的严格递减与唯一 $\gamma_0$; 两个端点值来自表
\ref{tab:envpoints}。其余各项由注 \ref{rem:env} 的有理包络方法逐项认证:
点值 (含 (iii)(iv) 的端点界) 见表 \ref{tab:envpoints}, 区间符号 (ii) 与
$Q_+<0$ 见表 \ref{tab:envsigns}, 小区间极值 $T_{A,B_2}\ge27/10$ 与
$T_C\ge19/10$ 见表 \ref{tab:envrange}, 全部单调性 (i)(iii)(iv)(v)(vi) 的导数
符号见表 \ref{tab:envderiv}。
\end{proof}

端点有理界 (26 项, 全部由注 \ref{rem:env} 的点值包络给出, 见表 \ref{tab:envpoints}; 用于表 \ref{tab:g1}):
\begin{equation}
	\begin{gathered}
	T_{A,B_2}(0.655)\ge\tfrac{11}{5},\ T_{A,B_2}(0.72)\ge\tfrac{13}{5},\ T_{A,B_2}(0.73)\ge\tfrac{13}{5},\\
	T_{A,B_2}(0.82)\ge2,\ T_{A,B_2}(0.83)\ge2,\ T_{A,B_2}(0.85)\ge\tfrac{19}{10},\ T_{A,B_2}(0.86)\ge\tfrac{47}{25},\\
	T_{A,M}(0.86)\ge\tfrac95,\ T_{A,M}(1.0014)\ge\tfrac35,\ T_{A,M}(1.0472)\ge\tfrac38,\\
	T_B(0.72)\ge\tfrac{3}{10},\ T_B(0.73)\ge\tfrac{3}{10},\ T_B(0.82)\ge\tfrac{3}{20},\ T_B(0.83)\ge\tfrac{3}{20},\\
	T_B(0.85)\ge\tfrac{1}{10},\ T_B(0.86)\ge\tfrac{1}{10},\ T_B(1.0014)\ge\tfrac{1}{25},\ T_B(1.0472)\ge\tfrac{1}{40},\\
	T_C(0.655)\ge\tfrac{57}{50},\ T_C(0.72)\ge\tfrac32,\ T_C(0.73)\ge\tfrac32,\ T_C(0.85)\ge\tfrac{19}{10},\\
	T_C(0.86)\ge\tfrac{19}{10},\ T_C(1.0014)\ge\tfrac43,\ T_C(1.0472)\ge\tfrac{11}{10},\ B_4(1.0472)\ge\tfrac{9}{25}.
	\end{gathered}
	\label{eq:endpoints}
\end{equation}

\begin{lemma}[二维界]\label{lem:j2bounds}
于 $[0.655,1.0472]\times[1,2]$ 上,
\begin{equation}
	\mathcal W_1+\mathcal W_2\le-T_A,\qquad
	\mathcal W_3\le-T_B,\qquad
	\mathcal W_4+\mathcal W_5\le-T_C,\qquad
	\mathcal W_6+\mathcal W_7+\mathcal W_8\le T_D.
	\label{eq:j2bounds}
\end{equation}
\end{lemma}

\begin{proof}
记 $z:=\mathrm{ct}^2$。由引理 \ref{lem:track}(i)(ii),
\begin{equation}
	\mathcal W_1+\mathcal W_2=z(1-z)\bigl(A^2\,\mathrm{cg}\,B_2-2A^3B_1z\bigr).
	\label{eq:w12}
\end{equation}
若 $\gamma\le\gamma_0$ (即 $B_1\ge0$), 括号 $\le A^2\,\mathrm{cg}\,B_2<0$
(引理 \ref{lem:brackets}(ii)), 故由 $z(1-z)\ge z_-(1-z_-)$ 与 $B_2<0$,
\begin{equation}
	\mathcal W_1+\mathcal W_2\le A^2\,\mathrm{cg}\,B_2\,z_-(1-z_-)
	=\frac{4A^2B_2\,\mathrm{sg}^2\,\mathrm{cg}^3}{D^4}
	\le\frac{4A^2B_2\,\mathrm{sg}^2\,\mathrm{cg}^4}{D^4}=-T_A,
\end{equation}
末步用 $0<\mathrm{cg}<1$, $B_2<0$。若 $\gamma\ge\gamma_0$ (即 $B_1\le0$),
用 $B_2=M+2A\,\mathrm{cg}\,B_1$ (引理 \ref{lem:j2dec}):
\begin{equation}
	A^2\,\mathrm{cg}\,B_2-2A^3B_1z
	=A^2\,\mathrm{cg}\,M+2A^3B_1(\mathrm{cg}^2-z)\le A^2\,\mathrm{cg}\,M<0,
\end{equation}
因 $z\le z_+=\mathrm{cg}^2$; 同一论证以 $M$ 代 $B_2$ 得
$\mathcal W_1+\mathcal W_2\le4A^2M\,\mathrm{sg}^2\,\mathrm{cg}^3/D^4
\le4A^2M\,\mathrm{sg}^2\,\mathrm{cg}^4/D^4=-T_A$。

对 $\mathcal W_3$: 由引理 \ref{lem:track}(iii),
\begin{equation}
	\mathcal W_3=-2A^3\,\mathrm{sg}\,(t\,\mathrm{st}\,\mathrm{ct}^5)
	\le-2A^3\,\mathrm{sg}\cdot\frac{2\tau\,\mathrm{sg}\,\mathrm{cg}^5}{D^6}
	=-\frac{4A^3\,\mathrm{sg}^2\,\tau\,\mathrm{cg}^5}{D^6}
	\le-\frac{2A^3\,\mathrm{sg}^2\,\tau\,\mathrm{cg}^5}{D^5}=-T_B,
\end{equation}
末步因 $D\le2$ (引理 \ref{lem:track}(v)) 使 $4/D^6\ge2/D^5$。

对 $\mathcal W_4+\mathcal W_5$: 用 $B_5=G_5+AB_4$,
\begin{equation}
	\mathcal W_4+\mathcal W_5
	=A\,\mathrm{sg}\,t\,\mathrm{st}\,\mathrm{ct}\bigl(AB_4\,\mathrm{ct}^2-\mathrm{cg}^2B_5\bigr)
	=A\,\mathrm{sg}\,t\,\mathrm{st}\,\mathrm{ct}\bigl(AB_4(\mathrm{ct}^2-\mathrm{cg}^2)-\mathrm{cg}^2G_5\bigr)
	\le-A\,\mathrm{sg}\,t\,\mathrm{st}\,\mathrm{ct}\,\mathrm{cg}^2G_5,
\end{equation}
因 $B_4>0$ 且 $\mathrm{ct}^2\le\mathrm{cg}^2$。又 $t\,\mathrm{st}\,\mathrm{ct}=\frac t2\sin2t\ge m$
(引理 \ref{lem:track}(iv)), 得 $\mathcal W_4+\mathcal W_5\le-mG_5A\,\mathrm{sg}\,\mathrm{cg}^2=-T_C$。

对 $\mathcal W_6+\mathcal W_7+\mathcal W_8$:
\begin{equation}
	\mathcal W_6+\mathcal W_7+\mathcal W_8
	=t^2\,\mathrm{cg}\,\mathrm{sg}^2\bigl(4A^2\,\mathrm{ct}^4-AB_7\,\mathrm{ct}^2
	+6\,\mathrm{cg}^2\,\mathrm{sg}^2\bigr)
	=t^2\,\mathrm{cg}\,\mathrm{sg}^2\,Q(z).
\end{equation}
$Q$ 关于 $z$ 凸 (二次项系数 $4A^2>0$), 故 $Q(z)\le\max(Q_-,Q_+)$; 由引理
\ref{lem:brackets}(iii), $Q_+<0$, 故 $Q(z)\le\max(Q_-,0)$。又 $t\le\tau$ 且
$\mathrm{cg}\,\mathrm{sg}^2>0$, 得
\begin{equation}
	\mathcal W_6+\mathcal W_7+\mathcal W_8
	\le\tau^2\,\mathrm{cg}\,\mathrm{sg}^2\max(Q_-,0)=T_D.
\end{equation}
\end{proof}

\begin{theorem}[$J_2^{(2)}$ 的完全解析上界]\label{thm:j2e1}
于 $[0.655,1.0472]\times[1,2]$ 上 $J_2^{(2)}(\gamma,q)<0$。
\end{theorem}

\begin{proof}
由引理 \ref{lem:j2dec} 与 \ref{lem:j2bounds},
\begin{equation}
	W=\sum_{j=1}^8\mathcal W_j\le-(T_A+T_B+T_C-T_D)=:-\mu.
\end{equation}
由 \eqref{eq:j2delta}--\eqref{eq:wdec}, $J_2^{(2)}=2A^2\,\mathrm{cg}\,W/\Delta^4$,
其中 $A>0$, $\mathrm{cg}>0$, $\Delta>0$; 故只需证 $\mu\ge139/100$。

端点有理界 (\eqref{eq:endpoints}) 全部由注 \ref{rem:env} 的点值包络给出
(表 \ref{tab:envpoints}), 单调性来自引理 \ref{lem:brackets}; 分段下界见表
\ref{tab:g1}。表中各行的组合方式如: 在
$[0.655,0.72]$ 上 $T_{A,B_2}$ 递增故 $T_A\ge T_{A,B_2}(0.655)\ge11/5$,
$T_B$ 递减故 $T_B\ge T_B(0.72)\ge3/10$, $T_C$ 递增故 $T_C\ge T_C(0.655)\ge57/50$;
在 $[0.72,0.724]$ 上 $T_B\ge T_B(0.73)\ge3/10$; 其余行同法。在 $[0.85,0.86]$ 上
$T_A\ge\min(T_{A,B_2}(0.86),T_{A,M}(0.86))\ge9/5$: 于 $[0.85,\gamma_0]$ 用
$T_{A,B_2}$ 递减且 $T_{A,B_2}(0.86)\ge47/25$, 于 $[\gamma_0,0.86]$ 用 $T_{A,M}$ 递减且
$T_{A,M}(0.86)\ge9/5$。在 $[0.655,1.0014]$ 上 $Q_-<0$ (因 $Q_-$ 递增且
$Q_-(1.0014)<0$), 故 $T_D=0$; 在 $[1.0014,1.0472]$ 上 $Q_-$ 递增且
$Q_-(1.0472)>0$, 故 $0\le\max(Q_-,0)\le Q_-(1.0472)\le33/200$, 且 $F$ 递增,
$F\le F(1.0472)\le63/100$, 于是
\begin{equation}
	T_D=F\max(Q_-,0)\le F(1.0472)\,Q_-(1.0472)
	\le\frac{63}{100}\cdot\frac{33}{200},
\end{equation}
而 $T_A,T_B,T_C$ 的递减性给出 $T_A\ge3/8$, $T_B\ge1/40$, $T_C\ge11/10$。故
\begin{equation}
	\mu\ge\frac38+\frac1{40}+\frac{11}{10}-\frac{63}{100}\cdot\frac{33}{200}
	=\frac{27921}{20000}>\frac{139}{100}.
\end{equation}
其余各段的 $\mu$ 值 (表 \ref{tab:g1} 末列) 均 $\ge139/100$。定理得证。
\end{proof}

\begin{table}[ht]
\centering
\caption{$\mu=T_A+T_B+T_C-T_D$ 的分段下界 (端点界见 \eqref{eq:endpoints})}
\label{tab:g1}
\begin{tabular}{lcccr}
\toprule
区间 & $T_A\ge$ & $T_B\ge$ & $T_C\ge$ & $\mu\ge$\\
\midrule
$[0.655,0.72]$ & $11/5$ & $3/10$ & $57/50$ & $91/25$\\
$[0.72,0.723]$ & $13/5$ & $3/10$ & $3/2$ & $22/5$\\
$[0.723,0.724]$ & $27/10$ & $3/10$ & $3/2$ & $9/2$\\
$[0.724,0.73]$ & $13/5$ & $3/10$ & $3/2$ & $22/5$\\
$[0.73,0.82]$ & $2$ & $3/20$ & $3/2$ & $73/20$\\
$[0.82,0.83]$ & $2$ & $3/20$ & $19/10$ & $81/20$\\
$[0.83,0.85]$ & $19/10$ & $1/10$ & $19/10$ & $39/10$\\
$[0.85,0.86]$ & $9/5$ & $1/10$ & $19/10$ & $19/5$\\
$[0.86,1.0014]$ & $3/5$ & $1/25$ & $4/3$ & $148/75$\\
$[1.0014,1.0472]$ & $3/8$ & $1/40$ & $11/10$ & $27921/20000$\\
\bottomrule
\end{tabular}
\end{table}

对第一相位与第二相位的 $q=1$ 线, 上面已给出完全初等的解析证明; 定理
\ref{thm:j2e1} 覆盖整个二维盒 $[0.655,1.0472]\times[1,2]$, 因而
\eqref{eq:Jcert} 不再依赖任何证书。


\begin{remark}[E3 侦察数据]\label{rem:explore}
以下范围来自高精度浮点扫描 (mpmath, 40 位), 只用于交叉检验与侦察, 不属于证明链。
$T_1$ 上: $G\in[-3.72,-2.81]$, $G_c\in[1.87,3.89]$,
$G_x\in[-1.07,1.28]$, $uG_x\in[-0.88,1.02]$, $x\Phi_q/(q+c\Phi_q)\in[0.68,0.84]$,
$J_1^{(2)}\in[9.01,18.61]$。$T_2$ 上: $G\in[-0.38,1.82]$,
$G_c\in[-2.64,0.26]$, $G_x\in[4.49,9.84]$, $x\Phi_q/(q+c\Phi_q)\in[1.40,1.86]$,
$J_2^{(2)}\in[-17.68,-5.99]$。这些数据与定理 \ref{thm:j1e1}、
\ref{thm:j2e1} 的解析结论一致; 旧引理 (M1$'$)--(M3$'$) 路线已被定理
\ref{thm:j2e1} 取代。
\end{remark}


\begin{remark}[$(LOG)$ 与 $(FP)$ 不等价]\label{rem:log}
伴随命题 \eqref{eq:LOG} 由定理 \ref{thm:LOG} 完全解析证明 (E1), 取代旧 128
叶盒证书路线 (第 \ref{sec:certs} 节)。这一对数比单调性与 $\widetilde F_e'<0$
逻辑上并不等价。本文唯一零点证明只使用独立证明的 $(FP)$ (即
\eqref{eq:keylemma}), 绝不从 $(LOG)$ 推导它。
\end{remark}

\subsection{对称线上恰有一个 good root}

\begin{theorem}[对称线 single crossing]\label{thm:single}
对每个 $R>1$, 函数 $S_R(\xi)$ 在 $(0,1/2)$ 中恰有一个零点 $\xi^*(R)$。并且
\begin{equation}
	S_R(\xi)<0\quad(0<\xi<\xi^*),\qquad
	S_R(\xi)>0\quad(\xi^*<\xi<1/2).
\end{equation}
该零点与对称 sign-consistent good root 双向等价。
\end{theorem}

\begin{proof}
由定理 \ref{thm:keylemma}, $\widetilde F_e$ 在 $(0,1/2)$ 严格递减。结合
\eqref{eq:Fepos} 与 \eqref{eq:Feneg}, 它在 $(0,1/2)$ 中恰有一个零点, 且在
$c\ge1/2$ 没有零点。映射 $\xi\mapsto c$ 严格递减, 而 \eqref{eq:dimred} 的比例
因子为正, 故得到所述 $S_R$ 符号顺序。

对称性使 $f_{\xi,1-\xi}(1-x)=f_{\xi,1-\xi}(x)$, 所以
\begin{equation}
	R_2(\xi,1-\xi)=R_1(\xi,1-\xi)=S_R(\xi).
\end{equation}
第一模态在内部为正, 第二模态关于中点为奇且在左半区间为正, 因此
\begin{equation}
	v(\xi)>0>v(1-\xi),\qquad 0<\xi<1/2.
\end{equation}
故 $S_R(\xi)=0$ 当且仅当 $(\xi,1-\xi)$ 是一个对称 good root。
\end{proof}

\section{O3a 的最终合成}

\begin{proof}[定理 \ref{thm:main} 的证明]
固定 $R>1$。定理 \ref{thm:single} 给出至少一个对称 good root
\begin{equation}
	(\xi^*(R),1-\xi^*(R)).
\end{equation}
另一方面, 任取参数三角形内的 good root $(a,b)$, 定理 \ref{thm:rigidity} 强制
$a+b=1$, 故它可写成
\begin{equation}
	(a,b)=(\xi,1-\xi),\qquad 0<\xi<1/2.
\end{equation}
由定理 \ref{thm:single} 的双向等价, $\xi$ 必是 $S_R$ 的唯一零点, 即
$\xi=\xi^*(R)$。因此整个参数三角形内恰有一个 good root。
\end{proof}

\begin{corollary}[canonical O3a]\label{cor:o3a}
若 $(a,b)$ 与 $(a',b')$ 是同一 $R>1$ 下的两个 sign-consistent good roots, 则
\begin{equation}
	(a,b)=(a',b').
\end{equation}
\end{corollary}

\section{边界、额外 sheets 与失败路线审计}

\begin{enumerate}
	\item \textbf{$R\to1^+$.} 定理对每个有限 $R>1$ 成立, 不需要随 $R\to1^+$
	保持一致的数值裕量。相位比证明甚至在形式极限 $m=1$ 时仍有 $\Psi'<0$。
	\item \textbf{$R\to\infty$.} 所有解析不等式对每个有限 $R$ 成立。第二模态外侧
	相位可非常接近 $\pi$, 但证明只需要严格小于 $\pi$, 不使用固定安全距离。
	\item \textbf{$a\to0^+$, $b\to1^-$, $a\to b$.} 目标定义在开参数三角形。实际
	good root 处被除的外侧正弦严格非零; 传输能量恒等式不除以中间相位
	$\sin\theta$, 所以即使 $b-a\to0$ 也无奇性。
	\item \textbf{多 residual sheets 与切触交点.} 对称刚性从任意 good root 出发,
	不需要选择``主分支''。因此额外 $R_2=0$ sheets、Jacobian 变号或切触均不能
	产生非对称 good root。
	\item \textbf{旧的 $h'>0$ 路线.} 独立高精度复算在 $R=1500$ 得到一个负导数值,
	故全局 $h'>0$ 命题为假。本文的相位比单调性是不同函数、不同变量上的严格
	结论, 不依赖该失败命题。
	\item \textbf{平方与归一化.} 传输步骤只使用正的平方振幅比, 且只作必要性
	推导; 没有从平方等式反推未验证的符号。所有 $n_k$ 严格为正, 并只在同一模态
	内消去。
\end{enumerate}

\section{证据层次与历史复现合同}\label{sec:certs}

本文当前版本\emph{不使用任何二维叶盒证书, 也不使用任何单变量事实验证器}。
早期版本中的四族证书——$J_1^{(2)}$ 的 16 叶盒、$J_2^{(2)}$ 的 67 叶盒、C4 区间段
$K>0$ 的 200 叶盒与伴随命题 $(LOG)$ 的 128 叶盒——已全部被纯解析证明 (E1)
取代并移除: $J_1^{(2)}$ 由定理 \ref{thm:j1e1}, $J_2^{(2)}$ 由定理 \ref{thm:j2e1},
$K>0$ 由引理 \ref{lem:corner} 的 C4 段, $(LOG)$ 由定理 \ref{thm:LOG}。主定理
\ref{thm:main} 的证明链只依赖 (E1) 解析证明; 其中 $J_2^{(2)}$ 所需的 55 项单变量
事实由注 \ref{rem:env} 的有理包络方法逐项给出 E1 证明, 证书总表见附录
\ref{app:cert}。

\subsection{历史证书一览 (已退役)}

旧 128 叶盒证书目标为 $(G_2-G_1)'<0$ 于闭盒 $Q$ (旧 $(LOG)$ 路线): 128 叶,
最接近零的认证上界 $-4.841603818885058$, 零失败。该路线已被定理 \ref{thm:LOG}
取代, 仅作历史记录; 原 16/67/200 叶盒族亦已分别随定理 \ref{thm:j1e1},
\ref{thm:j2e1} 与引理 \ref{lem:corner} 完成解析化而退役。每个叶盒族当时都用
有理端点验证了精确铺砌、无重叠和无遗漏 (叶面积总和等于盒面积), 并在独立重放
与二次引擎复算中全部通过; 这些历史记录不影响本文任何结论。

\subsection{单变量事实验证器 (历史, 已退役)}

更早版本对上述 55 项单变量事实使用十进制定向舍入区间运算引擎
(\texttt{misc/rigid\_dec.py}): $\pi$ 由 Machin 公式构造, $\sin,\cos$ 用带
显式交错级数余项界的 Taylor 展开, $\arctan$ 用交错级数或
$\pi/2-\arctan(1/x)$, 导数符号用对偶数自动微分加自适应细分, 端点比较用精确
有理分数; 引擎与 80 位 mpmath 做过 4800 次随机包含性检查, 零违反
(misc/\_test\_dec.py)。该引擎输出是``验证器认证的严格解析事实'', 不是
kernel-checked proof。自本版本起, 该引擎被注 \ref{rem:env} 的有理包络方法
(E1) 完全取代: 55 项事实全部改写为有限精确有理数不等式链 (附录 \ref{app:cert}
的证书总表), 不再依赖任何验证器内核。旧引擎三件套 (\texttt{misc/rigid\_dec.py},
\texttt{misc/zz\_verify\_e1\_dec.py}, \texttt{misc/e1\_facts\_ledger.json})
仅作历史复现记录, 不影响本文任何结论。

\subsection{内容哈希}

\begin{itemize}
	\item L5 (旧 $(LOG)$ 证书验证脚本, 已退役, 不再支撑任何结论):\\
	\texttt{\small 132e998f2a4f4807443c33e669435d6382de646b88be25d42e455251c7447f4a}
	\item L6 ($J_1^{(2)}$ E1 交叉检验脚本): \texttt{\small 64e24ace3117772b6cd2ea2ac53986a75cad6c3fd797b61369472ac87ec6ab04}
	\item L7--L9 (旧十进制区间验证器三件套 \texttt{rigid\_dec.py},\\
	\texttt{zz\_verify\_e1\_dec.py}, \texttt{e1\_facts\_ledger.json}, 已退役):
	\begin{quote}\ttfamily\scriptsize
	dd81278ed7a9e1ccf063cc446456c473546c1dbab73f6db7b60e11ec0d153525\\
	cad6c5ef56b7ccd38c2108def99916779cee77d387d8313c321912fdfed24bc4\\
	cc74fc5026866d33a367aad7dcd5152e6114751c9c2abe6a3b02e06b085cd9ef
	\end{quote}
	\item L10 (E1 有理包络证书生成器 \texttt{misc/e1\_certgen.py}):\\
	\texttt{\small 375209e2574aea15e3966b442316e2326070d75d4b9445d4bdb9ccf74dfec57c}
	\item L11 (E1 有理包络认证台账 \texttt{misc/e1\_cert\_ledger.json}):\\
	\texttt{\small ec9ce5ff7af7d9684bdd2097368e789e6f0b1dae798a04e62aef3d073fd68d30}
	\item L12 (E1 证书表生成器 \texttt{misc/e1\_cert\_tables.py}):\\
	\texttt{\small dce5c4538397257b823cd92cf1a7d4180a0ac24ba6e62b9d30d7d1efa33bb249}
\end{itemize}
L10--L12 为当前有理包络证书链的内容哈希; L5--L9 为历史产物, 仅作复现记录。
证书结论本身是附录 \ref{app:cert} 中的有限有理不等式链, 不依赖上述脚本的正确性。

\begin{remark}[可信度边界]\label{rem:trust}
本版本不依赖任何软件验证器内核。主定理 \ref{thm:main} 的证明链只由 (E1) 解析
证明组成: 闭式恒等式、初等不等式、单调性, 以及注 \ref{rem:env} 的有理包络方法
(附录 \ref{app:cert} 的证书总表——每一项都是一条由有限次精确有理数运算构成、
可人工复核的不等式链)。E3 扫描数据 (注 \ref{rem:explore} 等) 只用于交叉检验,
不构成结论依据。早期叶盒证书与十进制验证器均已退役, 仅作历史记录。
\end{remark}

\section{Blueprint 验收与证明谱系}

完整证明包的 SHA-256 为
\begin{center}
	\texttt{sha256:7e49f975a0dcd3cdc639515d5f2d72a004a4c36322991c067ce0f492c0a9f366}
\end{center}
最终不可变提案、校验、独立审稿和合并回执分别绑定为
\begin{center}
	\begin{tabular}{ll}
		\toprule
		提案 & \texttt{sha256:ed40b4b542dc02e8be2b2bfa208ef8ac3b2d9cc378599c73d7c81d432aca0495}\\
		校验 & \texttt{sha256:2cff40ebfaa6f951ea855e922381eda8618e22171daa9f3b1415bc48b534b38a}\\
		审稿 & \texttt{sha256:d5ac187991bd9805fae582adb4a66d685c8d15bb76e48954005dd94972503859}\\
		回执 & \texttt{sha256:87cb58e9d5d9d5582d98dcc61c0e6ffb6d5082ea853a44941335c9c5477a6cad}\\
		\bottomrule
	\end{tabular}
\end{center}
确定性接收器返回 \texttt{status=merged}。合并后的 canonical Blueprint SHA-256 为
\begin{center}
	\texttt{sha256:95cc5b2f7f8c23526642344a8fe1b6139611b73f679c3b90f2f37d9f654b303b}
\end{center}
合并后可信闭包包含 CLM-O3A, 使用的两个推理恰为 INF-O3A-PHASE-RIGIDITY 与
INF-O3A-FROM-SYMMETRY, 并且无矛盾。

\section{结论}

证明的真正新机制是``相位比刚性'': 中间区间的传输能量守恒把两个界面残差的共同
零点转化为同一严格单调函数在两个外侧相位上的等值。这个机制完全绕开了残差曲线
是否为单图、是否存在额外 sheet、以及分支差导数是否恒正等困难。它把二维全局
问题严格压缩到反射固定线; 对称线上由完全解析证明支持的 single-crossing 定理 (定理 \ref{thm:keylemma} 与 \ref{thm:single}) 随后完成
唯一性和存在性。

因此, 对每个 $R>1$, 全参数三角形中恰有一个 sign-consistent good root, 并且该根
必为对称根。

\appendix


\section{有理包络证书总表}\label{app:cert}

本节实现第 \ref{ss:j2e1} 子节注 \ref{rem:env} 所声明的全部证书。先给出两个
基础引理。

\begin{lemma}[交错级数包络]\label{lem:envseries}
对 $0<x<3/2$, 设
\begin{equation}
	S_m:=\sum_{k=0}^{m}\frac{(-1)^k x^{2k+1}}{(2k+1)!},\qquad
	C_m:=\sum_{k=0}^{m}\frac{(-1)^k x^{2k}}{(2k)!}\qquad(m\ge0)
\end{equation}
为 $\sin x$, $\cos x$ 的 Taylor 部分和 (各含 $m+1$ 项)。$\sin$ 的相邻项绝对值之比
$x^2/((2k+2)(2k+3))\le3/8<1$ 递减, 故部分和交替夹逼: 对 $m\ge0$,
\begin{equation}
	S_{2m}\ge\sin x\ge S_{2m+1},\qquad
	C_{2m}\ge\cos x\ge C_{2m+1}.
\end{equation}
($\cos$ 的 $m=0$ 情形 $C_0=1>\cos x>C_1=1-x^2/2$ 单独成立, 其后的项从 $k=2$ 起递
减, 故夹逼同样成立。) 由 Taylor 余项公式, 对一切 $x\ge0$ 更有
\begin{equation}
	|\sin x-S_m|\le\frac{x^{2m+3}}{(2m+3)!},\qquad
	|\cos x-C_m|\le\frac{x^{2m+2}}{(2m+2)!}.
\end{equation}
对 $0<x\le1$, $\arctan x=\sum_{k=0}^{\infty}(-1)^k x^{2k+1}/(2k+1)$ 同理 (部分和交替
夹逼, 余项 $\le x^{2m+3}/(2m+3)$); 对 $x>1$ 使用
$\arctan x=\pi/2-\arctan(1/x)$, $\pi$ 由 Machin 公式
$\pi=16\arctan(1/5)-4\arctan(1/239)$ 的两个反正切项包络组合得到有理界。
\end{lemma}

实际取值与包络宽度: 本文对 $\sin,\cos$ 取 12 项 ($m=11$), 在 $x\le3/2$ 时余项分别
$\le(3/2)^{25}/25!\approx1.6\times10^{-21}$ 与 $\le(3/2)^{24}/24!\approx2.7\times10^{-20}$,
均 $<10^{-12}$; 对 $\arctan$ 取 22 项 ($m=21$), 直接级数仅用于参数 $v\le1$, 余项
$\le v^{45}/45\le1/45$。Machin 参数 $v=1/5$, $1/239$ 的余项分别
$\le(1/5)^{45}/45\approx7.8\times10^{-34}$ 与 $\le(1/239)^{45}/45\approx2.1\times10^{-109}$,
故 $\pi$ 的认证宽度约 $2.5\times10^{-32}$。本文 $\tau(\gamma)=\arctan(2\tan\gamma)$ 在
$\gamma\in[0.655,1.0472]$ 恒有 $2\tan\gamma>1$ ($\tan\gamma\ge\tan(0.655)>0.655>1/2$),
故按 $\pi/2-\arctan(1/(2\tan\gamma))$ 计算, 其中参数 $v=1/(2\tan\gamma)$ 随 $\gamma$
递减, 在 $\gamma=0.655$ 处取最大值 $v\approx0.651$。因此表 \ref{tab:envprims} 的认证
宽度最坏为 $\tau(131/200)$ 行的约 $1.8\times10^{-10}$, 恰为上述 $\arctan$ 余项的二倍
$2\cdot v^{45}/45$; 其余原语 ($\sin,\cos,A,D$) 各行的宽度均 $\le10^{-23}$。全部证书的
最小裕量约 $2.6\times10^{-5}$ (表 \ref{tab:envpoints}), 故这些宽度远小于所有裕量,
不影响任何符号判定。

\begin{lemma}[泰勒模型判据]\label{lem:envtaylor}
设 $f$ 在 $[a,b]$ 上二阶连续可导, $c:=(a+b)/2$, $w:=(b-a)/2$。若 $f'(c)\in J$
(有理区间 $J=[j_-,j_+]$) 且 $\sup_{[a,b]}|f''|\le M$, 则对 $\gamma\in[a,b]$,
\begin{equation}
	f'(\gamma)\in[j_- - Mw,\; j_+ + Mw].
\end{equation}
特别地, $j_- - Mw>0$ 蕴含 $f'>0$ 于 $[a,b]$, $j_+ + Mw<0$ 蕴含 $f'<0$ 于
$[a,b]$。同理, 若 $f(c)\in J_0$ 且 $\sup_{[a,b]}|f'|\le M_0$, 则
\begin{equation}
	f(\gamma)\in J_0+[-M_0w,M_0w]\quad(\gamma\in[a,b]).
\end{equation}
\end{lemma}

\begin{proof}
$f'(\gamma)=f'(c)+f''(\xi)(\gamma-c)$ 且 $|\gamma-c|\le w$; 值版本用一阶 Taylor
余项 $f(\gamma)=f(c)+f'(\xi)(\gamma-c)$。
\end{proof}

注 \ref{rem:env}(b) 给出原语在小区间上的值域: $\sin\gamma$ 增, $\cos\gamma$ 减,
$A=\pi-\gamma$ 减, $D=\sqrt{1+3\sin^2\gamma}$ 增, $\tau=\arctan(2\tan\gamma)$ 增
($\tau'=2/D^2>0$); 故小区间端点处的点值包络直接给出小区间上的原语包络。导数
表达式由链式法则得到 ($\mathrm{d}A/\mathrm{d}\gamma=-1$,
$\mathrm{d}\sin\gamma/\mathrm{d}\gamma=\cos\gamma$,
$\mathrm{d}\cos\gamma/\mathrm{d}\gamma=-\sin\gamma$,
$\mathrm{d}D/\mathrm{d}\gamma=3\sin\gamma\cos\gamma/D$,
$\mathrm{d}\tau/\mathrm{d}\gamma=2/D^2$), 因此 $f'$ 与 $f''$ 在小区间上的包络
由同一原语包络经注 \ref{rem:env}(c) 的精确有理区间算术计算。以下四张表即由此
逐项生成; 每个区间都是一条有限精确有理数不等式链。生成脚本
misc/e1\_certgen.py 与认证台账 misc/e1\_cert\_ledger.json (内容哈希 L10--L12,
第 \ref{sec:certs} 节) 可复现全部数据, 但表中数据本身才是结论。

% ===== auto-generated certificate tables (e1_cert_tables.py) =====

\begin{table}[ht]
\centering
\footnotesize
\setlength{\tabcolsep}{3pt}
\caption{原语有理包络 (输入层): $\sin\gamma$, $\cos\gamma$, $\tau(\gamma)$, $A=\pi-\gamma$, $D=\sqrt{1+3\sin^2\gamma}$ 在 11 个主有理点处的认证区间。显示为 6 位小数向外取整 (显示区间包含认证区间); 认证宽度最坏为 $\tau(131/200)$ 行约 $1.8\times10^{-10}$ (22 项 $\arctan$ 级数在参数 $v=1/(2\tan\gamma)\approx0.651$ 处余项的二倍), 其余行宽度均 $\le10^{-23}$, 故多数行显示为单点。由注 \ref{rem:env} 与引理 \ref{lem:envseries} 计算。}\label{tab:envprims}
\begin{tabular}{lccccc}
\toprule
$\gamma$ & $\sin\gamma$ & $\cos\gamma$ & $\tau(\gamma)$ & $A$ & $D$\\
\midrule
$131/200$    & $0.609159\dots0.60916$     & $0.793047\dots0.793048$    & $0.993763\dots0.993764$    & $2.486592\dots2.486593$    & $1.453693\dots1.453694$    \\
$18/25$      & $0.659384\dots0.659385$    & $0.751805\dots0.751806$    & $1.052666\dots1.052667$    & $2.421592\dots2.421593$    & $1.518013\dots1.518014$    \\
$723/1000$   & $0.661637\dots0.661638$    & $0.749824\dots0.749825$    & $1.055265\dots1.055266$    & $2.418592\dots2.418593$    & $1.52095\dots1.520951$     \\
$181/250$    & $0.662386\dots0.662387$    & $0.749162\dots0.749163$    & $1.056129\dots1.05613$     & $2.417592\dots2.417593$    & $1.521929\dots1.52193$     \\
$73/100$     & $0.666869\dots0.66687$     & $0.745174\dots0.745175$    & $1.06129\dots1.061291$     & $2.411592\dots2.411593$    & $1.52779\dots1.527791$     \\
$41/50$      & $0.731145\dots0.731146$    & $0.682221\dots0.682222$    & $1.134271\dots1.134272$    & $2.321592\dots2.321593$    & $1.613605\dots1.613606$    \\
$83/100$     & $0.737931\dots0.737932$    & $0.674875\dots0.674876$    & $1.141908\dots1.141909$    & $2.311592\dots2.311593$    & $1.622845\dots1.622846$    \\
$17/20$      & $0.75128\dots0.751281$     & $0.659983\dots0.659984$    & $1.156927\dots1.156928$    & $2.291592\dots2.291593$    & $1.641117\dots1.641118$    \\
$43/50$      & $0.757842\dots0.757843$    & $0.652437\dots0.652438$    & $1.164312\dots1.164313$    & $2.281592\dots2.281593$    & $1.650144\dots1.650145$    \\
$5007/5000$  & $0.842226\dots0.842227$    & $0.539123\dots0.539124$    & $1.26104\dots1.261041$     & $2.140192\dots2.140193$    & $1.768625\dots1.768626$    \\
$1309/1250$  & $0.866026\dots0.866027$    & $0.499997\dots0.499998$    & $1.289762\dots1.289763$    & $2.094392\dots2.094393$    & $1.802777\dots1.802778$    \\
\bottomrule
\end{tabular}
\end{table}

\begin{table}[ht]
\centering
\small
\caption{点值证书 (输出层): 每个事实的认证区间 (12 位小数向外取整, 区间宽度 $\le10^{-12}$ 故通常显示为单点) 与所需有理界; 裕量 = 认证界到目标界的距离, 严格为正。}\label{tab:envpoints}
\begin{tabular}{lll}
\toprule
事实 & 认证值 & 裕量\\
\midrule
$B_1(0.85)\ge1/200$ & $0.009851718323\;\ge\;1/200$ & $4.9e-03$\\
$B_1(0.86)\le-1/50$ & $-0.027088591501\;\le\;-1/50$ & $7.1e-03$\\
$B_4(1.0472)\ge9/25$ & $0.362311210636\;\ge\;9/25$ & $2.3e-03$\\
$Q_-(1.0014)\le-1/10000$ & $-0.000163114796\;\le\;-1/10000$ & $6.3e-05$\\
$Q_-(1.0472)\le33/200$ & $0.16460378084\;\le\;33/200$ & $4.0e-04$\\
$F(1.0472)\le63/100$ & $0.623807276062\;\le\;63/100$ & $6.2e-03$\\
$\tau(1.0472)<13/10$ & $1.289762932247\;\le\;13/10$ & $1.0e-02$\\
$T_{A,B_2}(0.655)\ge11/5$ & $2.289785651689\;\ge\;11/5$ & $9.0e-02$\\
$T_{A,B_2}(0.72)\ge13/5$ & $2.701505973438\;\ge\;13/5$ & $1.0e-01$\\
$T_{A,B_2}(0.73)\ge13/5$ & $2.699537088766\;\ge\;13/5$ & $1.0e-01$\\
$T_{A,B_2}(0.82)\ge2$ & $2.225883001375\;\ge\;2$ & $2.3e-01$\\
$T_{A,B_2}(0.83)\ge2$ & $2.142798843267\;\ge\;2$ & $1.4e-01$\\
$T_{A,B_2}(0.85)\ge19/10$ & $1.96925635975\;\ge\;19/10$ & $6.9e-02$\\
$T_{A,B_2}(0.86)\ge47/25$ & $1.880222182803\;\ge\;47/25$ & $2.2e-04$\\
$T_{A,M}(0.86)\ge9/5$ & $1.856652401546\;\ge\;9/5$ & $5.7e-02$\\
$T_{A,M}(1.0014)\ge3/5$ & $0.6099754818\;\ge\;3/5$ & $1.0e-02$\\
$T_{A,M}(1.0472)\ge3/8$ & $0.385308867365\;\ge\;3/8$ & $1.0e-02$\\
$T_B(0.72)\ge3/10$ & $0.387304002615\;\ge\;3/10$ & $8.7e-02$\\
$T_B(0.73)\ge3/10$ & $0.365449806165\;\ge\;3/10$ & $6.5e-02$\\
$T_B(0.82)\ge3/20$ & $0.204999224386\;\ge\;3/20$ & $5.5e-02$\\
$T_B(0.83)\ge3/20$ & $0.19105580724\;\ge\;3/20$ & $4.1e-02$\\
$T_B(0.85)\ge1/10$ & $0.165317280626\;\ge\;1/10$ & $6.5e-02$\\
$T_B(0.86)\ge1/10$ & $0.153481770234\;\ge\;1/10$ & $5.3e-02$\\
$T_B(1.0014)\ge1/25$ & $0.046157042222\;\ge\;1/25$ & $6.2e-03$\\
$T_B(1.0472)\ge1/40$ & $0.029168022035\;\ge\;1/40$ & $4.2e-03$\\
$T_C(0.655)\ge57/50$ & $1.145710285107\;\ge\;57/50$ & $5.7e-03$\\
$T_C(0.72)\ge3/2$ & $1.665131702973\;\ge\;3/2$ & $1.7e-01$\\
$T_C(0.73)\ge3/2$ & $1.721566781687\;\ge\;3/2$ & $2.2e-01$\\
$T_C(0.85)\ge19/10$ & $1.942123508194\;\ge\;19/10$ & $4.2e-02$\\
$T_C(0.86)\ge19/10$ & $1.926812153114\;\ge\;19/10$ & $2.7e-02$\\
$T_C(1.0014)\ge4/3$ & $1.360828709937\;\ge\;4/3$ & $2.7e-02$\\
$T_C(1.0472)\ge11/10$ & $1.105239258133\;\ge\;11/10$ & $5.2e-03$\\
$h(0.655)\ge m$ & $0.316425571654\;\ge\;791/2500$ & $2.6e-05$\\
$h(13/10)\ge m$ & $0.335075891684\;\ge\;791/2500$ & $1.9e-02$\\
\bottomrule
\end{tabular}
\end{table}

\begin{table}[ht]
\centering
\small
\caption{区间符号证书: 在小区间上量值的认证区间 (值泰勒模型, 注 \ref{rem:env}); 每个小区间上 $f$ 的认证上界 $<0$ (或下界 $>0$), 裕量为该界到零的距离。数值显示为 6 位小数向外取整 (显示区间包含认证区间)。}\label{tab:envsigns}
\begin{tabular}{llll}
\toprule
事实 & 小区间 & 认证值区间 & 裕量\\
\midrule
$B_2<0$ & $[0.655,\,0.75305]$ & $[-5.621046,\,-2.481017]$ & $2.48e+00$\\
 & $[0.75305,\,0.8511]$ & $[-6.910167,\,-4.731748]$ & $4.73e+00$\\
 & $[0.8511,\,0.94915]$ & $[-7.29394,\,-6.056586]$ & $6.06e+00$\\
 & $[0.94915,\,1.0472]$ & $[-7.237108,\,-6.110895]$ & $6.11e+00$\\
$M<0$ & $[0.655,\,0.8511]$ & $[-7.74232,\,-4.976264]$ & $4.98e+00$\\
 & $[0.8511,\,1.0472]$ & $[-7.538029,\,-4.231683]$ & $4.23e+00$\\
$B_4>0$ & $[0.655,\,0.704025]$ & $[6.876802,\,8.125115]$ & $6.88e+00$\\
 & $[0.704025,\,0.75305]$ & $[5.744307,\,6.950459]$ & $5.74e+00$\\
 & $[0.75305,\,0.802075]$ & $[4.6636,\,5.81744]$ & $4.66e+00$\\
 & $[0.802075,\,0.8511]$ & $[3.643712,\,4.735952]$ & $3.64e+00$\\
 & $[0.8511,\,0.900125]$ & $[2.692724,\,3.715025]$ & $2.69e+00$\\
 & $[0.900125,\,0.94915]$ & $[1.817684,\,2.762739]$ & $1.82e+00$\\
 & $[0.94915,\,0.998175]$ & $[1.024557,\,1.886149]$ & $1.02e+00$\\
 & $[0.998175,\,1.0472]$ & $[0.318167,\,1.091221]$ & $3.18e-01$\\
$G_5>0$ & $[0.655,\,0.704025]$ & $[3.627494,\,5.639475]$ & $3.63e+00$\\
 & $[0.704025,\,0.75305]$ & $[5.215643,\,6.895374]$ & $5.22e+00$\\
 & $[0.75305,\,0.802075]$ & $[6.479632,\,7.82551]$ & $6.48e+00$\\
 & $[0.802075,\,0.8511]$ & $[7.419682,\,8.434798]$ & $7.42e+00$\\
 & $[0.8511,\,0.900125]$ & $[8.040328,\,8.732598]$ & $8.04e+00$\\
 & $[0.900125,\,0.94915]$ & $[8.350901,\,8.731763]$ & $8.35e+00$\\
 & $[0.94915,\,0.998175]$ & $[8.094607,\,8.718978]$ & $8.09e+00$\\
 & $[0.998175,\,1.0472]$ & $[7.567267,\,8.435353]$ & $7.57e+00$\\
$Q_+<0$ & $[0.655,\,0.704025]$ & $[-3.884826,\,-3.357854]$ & $3.36e+00$\\
 & $[0.704025,\,0.75305]$ & $[-3.799617,\,-3.181378]$ & $3.18e+00$\\
 & $[0.75305,\,0.802075]$ & $[-3.598361,\,-2.916551]$ & $2.92e+00$\\
 & $[0.802075,\,0.8511]$ & $[-3.305858,\,-2.588393]$ & $2.59e+00$\\
 & $[0.8511,\,0.900125]$ & $[-2.94777,\,-2.221567]$ & $2.22e+00$\\
 & $[0.900125,\,0.94915]$ & $[-2.549446,\,-1.839318]$ & $1.84e+00$\\
 & $[0.94915,\,0.998175]$ & $[-2.134834,\,-1.462534]$ & $1.46e+00$\\
 & $[0.998175,\,1.0472]$ & $[-1.725516,\,-1.108973]$ & $1.11e+00$\\
\bottomrule
\end{tabular}
\end{table}

\begin{table}[ht]
\centering
\small
\caption{小区间极值证书: 两个``区间下界''事实, 由值泰勒模型直接给出 (不依赖单调性); 每小区间认证下界与目标界之差 (裕量) 严格为正。数值显示为 6 位小数向外取整 (显示区间包含认证区间)。}\label{tab:envrange}
\begin{tabular}{llll}
\toprule
事实 & 小区间 & 认证值区间 & 裕量\\
\midrule
$T_{A,B_2}\ge27/10$ 于 $[0.723,0.724]$ & $[0.723,\,0.72325]$ & $[2.702392,\,2.702434]$ & $2.39e-03$\\
 & $[0.72325,\,0.7235]$ & $[2.702408,\,2.702441]$ & $2.41e-03$\\
 & $[0.7235,\,0.72375]$ & $[2.702413,\,2.702441]$ & $2.41e-03$\\
 & $[0.72375,\,0.724]$ & $[2.702402,\,2.70244]$ & $2.40e-03$\\
$T_C\ge19/10$ 于 $[0.82,0.83]$ & $[0.82,\,0.8225]$ & $[1.959822,\,1.960719]$ & $5.98e-02$\\
 & $[0.8225,\,0.825]$ & $[1.959879,\,1.960715]$ & $5.99e-02$\\
 & $[0.825,\,0.8275]$ & $[1.959435,\,1.960581]$ & $5.94e-02$\\
 & $[0.8275,\,0.83]$ & $[1.958681,\,1.960134]$ & $5.87e-02$\\
\bottomrule
\end{tabular}
\end{table}

{\setlength{\tabcolsep}{3pt}
\begin{longtable}{llll}
\caption{导数符号证书: 每个小区间 $[a,b]$ 上 $f'$ 的认证区间 (泰勒模型 $f'(\gamma)\in f'(c)+f''(\cdot)[-w,w]$, $w=(b-a)/2$, 注 \ref{rem:env}) 与裕量。所有 $f'$ 认证区间与零严格分离, 故各单调性成立。数值显示为 6 位小数向外取整 (显示区间包含认证区间)。}\label{tab:envderiv}\\
\toprule
事实 & 小区间 & $f'$ 认证区间 & 裕量\\
\midrule
\endfirsthead
\toprule
事实 & 小区间 & $f'$ 认证区间 & 裕量\\
\midrule
\endhead
\bottomrule
\endlastfoot
$Q_-$ 递增 & $[0.655,\,0.75305]$ & $[7.553438,\,19.505772]$ & $7.55e+00$\\
 & $[0.75305,\,0.8511]$ & $[6.085026,\,16.088964]$ & $6.09e+00$\\
 & $[0.8511,\,0.94915]$ & $[3.675697,\,11.739948]$ & $3.68e+00$\\
 & $[0.94915,\,1.0472]$ & $[1.316453,\,7.460499]$ & $1.32e+00$\\
$F$ 递增于 $[1.0014,1.0472]$ & $[1.0014,\,1.0243]$ & $[0.325025,\,0.464678]$ & $3.25e-01$\\
 & $[1.0243,\,1.0472]$ & $[0.216445,\,0.363047]$ & $2.16e-01$\\
$T_{A,B_2}$ 递增于 $[0.655,0.72]$ & $[0.655,\,0.663125]$ & $[10.989882,\,13.431781]$ & $1.10e+01$\\
 & $[0.663125,\,0.67125]$ & $[9.172959,\,11.495877]$ & $9.17e+00$\\
 & $[0.67125,\,0.679375]$ & $[7.454505,\,9.658438]$ & $7.45e+00$\\
 & $[0.679375,\,0.6875]$ & $[5.834324,\,7.919797]$ & $5.83e+00$\\
 & $[0.6875,\,0.695625]$ & $[4.311774,\,6.279795]$ & $4.31e+00$\\
 & $[0.695625,\,0.70375]$ & $[2.885816,\,4.737828]$ & $2.89e+00$\\
 & $[0.70375,\,0.711875]$ & $[1.555048,\,3.292885]$ & $1.56e+00$\\
 & $[0.711875,\,0.72]$ & $[0.317746,\,1.943591]$ & $3.18e-01$\\
$T_{A,B_2}$ 递增于 $[0.72,0.723]$ & $[0.72,\,0.7215]$ & $[0.292388,\,0.524577]$ & $2.92e-01$\\
 & $[0.7215,\,0.723]$ & $[0.075482,\,0.304453]$ & $7.55e-02$\\
$T_{A,B_2}$ 递减于 $[0.724,0.73]$ & $[0.724,\,0.727]$ & $[-0.506868,\,-0.038831]$ & $3.88e-02$\\
 & $[0.727,\,0.73]$ & $[-0.91473,\,-0.459778]$ & $4.60e-01$\\
$T_{A,B_2}$ 递减于 $[0.73,0.85]$ & $[0.73,\,0.745]$ & $[-3.377906,\,-0.337881]$ & $3.38e-01$\\
 & $[0.745,\,0.76]$ & $[-4.898371,\,-2.251203]$ & $2.25e+00$\\
 & $[0.76,\,0.775]$ & $[-6.154326,\,-3.873948]$ & $3.87e+00$\\
 & $[0.775,\,0.79]$ & $[-7.172924,\,-5.216718]$ & $5.22e+00$\\
 & $[0.79,\,0.805]$ & $[-7.964226,\,-6.309294]$ & $6.31e+00$\\
 & $[0.805,\,0.82]$ & $[-8.54913,\,-7.171784]$ & $7.17e+00$\\
 & $[0.82,\,0.835]$ & $[-8.949286,\,-7.823922]$ & $7.82e+00$\\
 & $[0.835,\,0.85]$ & $[-9.184482,\,-8.287013]$ & $8.29e+00$\\
$T_{A,B_2}$ 递减于 $[0.85,0.86]$ & $[0.85,\,0.855]$ & $[-8.940649,\,-8.819617]$ & $8.82e+00$\\
 & $[0.855,\,0.86]$ & $[-8.979665,\,-8.876342]$ & $8.88e+00$\\
$T_{A,M}$ 递减于 $[0.85,0.86]$ & $[0.85,\,0.855]$ & $[-12.415415,\,-12.162098]$ & $1.22e+01$\\
 & $[0.855,\,0.86]$ & $[-12.20846,\,-11.952722]$ & $1.20e+01$\\
$T_{A,M}$ 递减于 $[0.86,1.0472]$ & $[0.86,\,0.9068]$ & $[-14.110062,\,-7.790287]$ & $7.79e+00$\\
 & $[0.9068,\,0.9536]$ & $[-11.487878,\,-6.135539]$ & $6.14e+00$\\
 & $[0.9536,\,1.0004]$ & $[-8.929862,\,-4.563717]$ & $4.56e+00$\\
 & $[1.0004,\,1.0472]$ & $[-6.626245,\,-3.195396]$ & $3.20e+00$\\
$T_B$ 递减 & $[0.655,\,0.704025]$ & $[-3.41569,\,-1.652157]$ & $1.65e+00$\\
 & $[0.704025,\,0.75305]$ & $[-2.877123,\,-1.436773]$ & $1.44e+00$\\
 & $[0.75305,\,0.802075]$ & $[-2.340306,\,-1.17864]$ & $1.18e+00$\\
 & $[0.802075,\,0.8511]$ & $[-1.839316,\,-0.926014]$ & $9.26e-01$\\
 & $[0.8511,\,0.900125]$ & $[-1.399023,\,-0.700365]$ & $7.00e-01$\\
 & $[0.900125,\,0.94915]$ & $[-1.031144,\,-0.510113]$ & $5.10e-01$\\
 & $[0.94915,\,0.998175]$ & $[-0.73638,\,-0.357204]$ & $3.57e-01$\\
 & $[0.998175,\,1.0472]$ & $[-0.508751,\,-0.239663]$ & $2.40e-01$\\
$T_C$ 递增于 $[0.655,0.82]$ & $[0.655,\,0.665312\dots0.665313]$ & $[9.329366,\,10.048058]$ & $9.33e+00$\\
 & $[0.665312\dots0.665313,\,0.675625]$ & $[8.69327,\,9.420094]$ & $8.69e+00$\\
 & $[0.675625,\,0.685937\dots0.685938]$ & $[8.050772,\,8.783253]$ & $8.05e+00$\\
 & $[0.685937\dots0.685938,\,0.69625]$ & $[7.40424,\,8.139955]$ & $7.40e+00$\\
 & $[0.69625,\,0.706562\dots0.706563]$ & $[6.755983,\,7.492568]$ & $6.76e+00$\\
 & $[0.706562\dots0.706563,\,0.716875]$ & $[6.108247,\,6.843404]$ & $6.11e+00$\\
 & $[0.716875,\,0.727187\dots0.727188]$ & $[5.46321,\,6.194709]$ & $5.46e+00$\\
 & $[0.727187\dots0.727188,\,0.7375]$ & $[4.822972,\,5.548662]$ & $4.82e+00$\\
 & $[0.7375,\,0.747812\dots0.747813]$ & $[4.189558,\,4.907367]$ & $4.19e+00$\\
 & $[0.747812\dots0.747813,\,0.758125]$ & $[3.564905,\,4.272849]$ & $3.56e+00$\\
 & $[0.758125,\,0.768437\dots0.768438]$ & $[2.950864,\,3.647048]$ & $2.95e+00$\\
 & $[0.768437\dots0.768438,\,0.77875]$ & $[2.349191,\,3.031815]$ & $2.35e+00$\\
 & $[0.77875,\,0.789062\dots0.789063]$ & $[1.761548,\,2.42891]$ & $1.76e+00$\\
 & $[0.789062\dots0.789063,\,0.799375]$ & $[1.189498,\,1.839997]$ & $1.19e+00$\\
 & $[0.799375,\,0.809687\dots0.809688]$ & $[0.6345,\,1.266639]$ & $6.35e-01$\\
 & $[0.809687\dots0.809688,\,0.82]$ & $[0.097913,\,0.7103]$ & $9.79e-02$\\
$T_C$ 递减于 $[0.83,1.0472]$ & $[0.83,\,0.85715]$ & $[-1.89793,\,-0.130822]$ & $1.31e-01$\\
 & $[0.85715,\,0.8843]$ & $[-2.987807,\,-1.408223]$ & $1.41e+00$\\
 & $[0.8843,\,0.91145]$ & $[-3.904116,\,-2.524433]$ & $2.52e+00$\\
 & $[0.91145,\,0.9386]$ & $[-4.639826,\,-3.467611]$ & $3.47e+00$\\
 & $[0.9386,\,0.96575]$ & $[-5.194292,\,-4.229766]$ & $4.23e+00$\\
 & $[0.96575,\,0.9929]$ & $[-5.569943,\,-4.80959]$ & $4.81e+00$\\
 & $[0.9929,\,1.02005]$ & $[-5.774352,\,-5.209708]$ & $5.21e+00$\\
 & $[1.02005,\,1.0472]$ & $[-5.81838,\,-5.437675]$ & $5.44e+00$\\
\end{longtable}}

\section{符号速查}

\begin{center}
\begin{tabular}{ll}
\toprule
符号 & 含义\\
\midrule
$R=q^2=m^2$ & 高密度区间上的权值, 高于背景值 $1$\\
$(a,b)$ & 一般势垒界面参数\\
$\xi$ & 对称线上左界面, $(a,b)=(\xi,1-\xi)$\\
$s_k=\sqrt{\lambda_k}$ & 第 $k$ 个谱频率\\
$\alpha,\beta$ & 一般 good root 的左右第一模态外侧相位\\
$\tau=s_2/s_1$ & 前两频率之比\\
$\alpha_k=s_k\xi$ & 对称半问题中的第 $k$ 个左相位\\
$c=q(1/2-\xi)/\xi$ & 对称几何的无量纲坐标\\
$\Phi_q(x)$ & $\cos^2x+q^2\sin^2x$\\
$\widetilde M_f,\widetilde F_e$ & 对称谱隙导数的单变量降维量\\
$S_R(\xi)$ & canonical 对称界面残差\\
$G_k$ & 沿第 $k$ 条相位曲线的 $\log\widetilde M_{fk}$ 导数\\
\bottomrule
\end{tabular}
\end{center}

\section{可复现产物位置}

本文依据以下独立项目产物整理:
\begin{itemize}
	\item 完整证明包: \texttt{runs/R-20260808T143337Z-o3a-c1/artifacts/proof-INF-O3A-COMPLETE.md};
	\item 相位比解析推导: \texttt{routes/analytic\_algebra/derivations.md};
	\item 机制不同的对抗审计: \texttt{routes/adversarial\_compute/symmetry\_chain\_audit.md};
	\item 对称线独立审计: \texttt{routes/symmetric\_line\_audit/independent\_audit.md};
	\item 证书重放输出: \texttt{routes/symmetric\_line\_audit/replay/replay\_output.txt};
	\item canonical 合并回执: \texttt{statistics/submissions/SUB-20260809-0003-O3APROOF2/receipt.json}。
\end{itemize}

本仓库内的审计脚本 (2026-08-10 复跑, 全部通过):
\begin{itemize}
	\item \texttt{scripts/audit\_o3a\_pdf\_part1.py};
	\item \texttt{scripts/audit\_o3a\_pdf\_part2.py};
	\item \texttt{scripts/audit\_o3a\_pdf\_part2b.py};
	\item \texttt{scripts/audit\_o3a\_pdf\_part2c.py};
	\item \texttt{scripts/audit\_o3a\_pdf\_part3.py};
	\item \texttt{scripts/audit\_o3a\_pdf\_part4.py};
	\item 大 $R$ 精确定位: \texttt{scripts/\_audit\_cstar.py};
	\item 大 $R$ 残差复核: \texttt{scripts/\_tmp\_verify\_r1e6.py};
\end{itemize}
2026-08-10 复跑说明: 上述 8 个脚本全部通过; 其中 \texttt{part2b} (特征值网格上界取 $3\pi$, $R$ 列表不含 $10^6$) 与 \texttt{part2c} ($\xi$ 扫描加密至 $0.4999995$, mpmath 精修全程保持 \texttt{mpf}) 为本次网格/精度修正后的版本. 以上均为 E3 交叉检验, 不构成任何 E1 前提.
定理 \ref{thm:j2e1} 的 E1 验证产物:
\begin{itemize}
	\item 有理包络证书链: \texttt{misc/e1\_certgen.py} (生成器),
	\texttt{misc/e1\_cert\_ledger.json} (认证台账), \texttt{misc/e1\_cert\_tables.py}
	(证书表生成器); 内容哈希 L10--L12, 证书总表见附录 \ref{app:cert};
	\item 精确有理区间内核: \texttt{misc/rigid1d.py} (交错级数包络 + Machin $\pi$;
	2026-08-09 修复 \texttt{I.sqrt} 的 $+1$ 单位错误);
	\item 代数分解精确验证: \texttt{misc/\_verify\_identity.py},
	\texttt{misc/zz\_rebuild\_check1.py}, 分子数据文件 \texttt{misc/t3\_NJ2.json};
	\item 旧十进制区间验证器三件套 (\texttt{misc/rigid\_dec.py},\\
	\texttt{misc/zz\_verify\_e1\_dec.py}, \texttt{misc/e1\_facts\_ledger.json})
	已退役, 仅作历史复现 (内容哈希 L7--L9)。
\end{itemize}

\end{document}
